% ============================================================================
%  profiler.tex -- Reference manual for profiler.py, the SM8850P
%  heterogeneous profiler wrapper API.
%
%  Compile with:  pdflatex profiler.tex   (twice, for the ToC / refs)
%  Only standard TeX Live / Overleaf packages are used; no shell-escape.
%
%  Revision 2.0 supersedes the SM8750 edition. The board changed
%  (SM8750 -> SM8850P) and the API grew: hardware counters, threading,
%  power estimation, converted networks, custom kernels, and eight more
%  workloads. See docs/board-migration-sm8750-to-sm8850.md.
% ============================================================================
\documentclass[11pt,a4paper]{article}

\usepackage[T1]{fontenc}
\usepackage[utf8]{inputenc}
\usepackage{lmodern}
\usepackage{textcomp}
\usepackage{amsmath}
\usepackage[margin=1in]{geometry}
\usepackage{microtype}
\usepackage{booktabs}
\usepackage{longtable}
\usepackage{array}
\usepackage{xcolor}
\usepackage{listings}
\usepackage{enumitem}
\usepackage{fancyhdr}
\usepackage[hidelinks]{hyperref}

% ---------------------------------------------------------------- colours ---
\definecolor{codebg}{RGB}{247,247,249}
\definecolor{coderule}{RGB}{220,220,226}
\definecolor{kw}{RGB}{0,90,160}
\definecolor{cmt}{RGB}{110,120,130}
\definecolor{str}{RGB}{150,60,40}

% --------------------------------------------------------------- listings ---
\lstdefinestyle{py}{
  language=Python,
  basicstyle=\ttfamily\footnotesize,
  keywordstyle=\color{kw}\bfseries,
  commentstyle=\color{cmt}\itshape,
  stringstyle=\color{str},
  backgroundcolor=\color{codebg},
  frame=single,
  rulecolor=\color{coderule},
  framesep=5pt,
  showstringspaces=false,
  columns=fullflexible,
  keepspaces=true,
  breaklines=true,
  breakatwhitespace=true,
  xleftmargin=0pt,
  aboveskip=8pt,
  belowskip=8pt,
}
\lstdefinestyle{sh}{
  basicstyle=\ttfamily\footnotesize,
  backgroundcolor=\color{codebg},
  frame=single,
  rulecolor=\color{coderule},
  framesep=5pt,
  showstringspaces=false,
  columns=fullflexible,
  keepspaces=true,
  breaklines=true,
  aboveskip=8pt,
  belowskip=8pt,
}
\lstdefinestyle{plain}{
  basicstyle=\ttfamily\scriptsize,
  backgroundcolor=\color{white},
  frame=none,
  columns=fullflexible,
  keepspaces=true,
  aboveskip=6pt,
  belowskip=6pt,
}
\lstset{style=py}

% ---------------------------------------------------------------- helpers ---
% Long identifiers may break after an underscore instead of running into the margin.
\newcommand{\code}[1]{\texttt{\let\_\breakableunderscore #1}}
\newcommand{\breakableunderscore}{\textunderscore\allowbreak}
\setlength{\emergencystretch}{3em}
\newcommand{\us}{$\mu$s}
\newcommand{\arrow}{$\rightarrow$}
\newcommand{\apient}[1]{\subsection{\texttt{#1}}}

% ----------------------------------------------------------------- header ---
\pagestyle{fancy}
\fancyhf{}
\lhead{\small SM8850P Profiler Wrapper API}
\rhead{\small \thepage}
\renewcommand{\headrulewidth}{0.4pt}

\title{\vspace{-1.5cm}\textbf{The SM8850P Profiler Wrapper API}\\[0.3cm]
       \large A dispatch, counter and measurement interface for the four compute
       engines\\ of the Qualcomm Snapdragon 8 Elite Gen 5 (SM8850P)}
\author{Snapdragon Profiler project --- \code{profiler.py}}
\date{Document revision 2.0}

\begin{document}
\maketitle
\thispagestyle{fancy}

\begin{abstract}
\noindent
\code{profiler.py} is a single-file, standard-library-only Python API that
dispatches benchmark kernels to a Qualcomm HDK8850 board over \code{adb} and
returns parsed, comparable measurements. It exposes twelve workloads across
four compute engines (CPU, GPU, DSP, NPU) behind one uniform call; collects
ARMv9 hardware performance counters alongside any CPU dispatch; runs converted
neural networks through QNN; estimates total-board power from a calibrated
model; and supports serial execution, genuine hardware-parallel execution, and
dependency-ordered DAG execution.

This document is the complete reference: every class, every argument, every
field, and the measurement caveats that decide whether a number you get back is
trustworthy. It is written to be read in order once and then used as a lookup.

\medskip
\noindent\textbf{One standing warning about catalogues.} The kernel and
parameter tables reproduced here are hand-maintained and can lag the code.
\code{python profiler.py -{}-catalogue} and \code{python profiler.py
-{}-list-counters} are \emph{generated from the registry at runtime} and are
therefore always current. Where this document and those commands disagree,
the commands are right.
\end{abstract}

\tableofcontents
\newpage

% ============================================================================
\section{Overview}
% ============================================================================

\subsection{What this is}

The Snapdragon Profiler project builds and characterises benchmark kernels for
the SM8850P SoC. Underneath, every measurement is an on-device binary invoked
over \code{adb} with a backend-specific working directory, library path,
affinity mask, privilege wrapper and flag set, whose stdout is then scraped for
timing fields that differ per backend.

\code{profiler.py} is the API over that machinery. It turns

\begin{lstlisting}[style=sh]
adb shell "cd /data/local/tmp/profiler-kernels/npu && LD_LIBRARY_PATH=... \
  ADSP_LIBRARY_PATH=... ./conv_qnn_fp16 --perf-mode sustained --reps 5 --cin 8 --cout 8"
\end{lstlisting}

\noindent into

\begin{lstlisting}
p.run(convolution(Backend.NPU, cin=8, cout=8))
\end{lstlisting}

\noindent and returns a \code{Result} whose \code{compute\_us}, \code{gflops}
and \code{gbps} fields mean the same thing regardless of which engine produced
them.

\subsection{The hardware it targets}
\label{sec:hardware}

The course hardware is the \textbf{HDK8850}, and several numbers in the API
are properties of that part rather than of the software.

\begin{longtable}{@{}p{2.6cm}p{12.2cm}@{}}
\toprule
\textbf{Unit} & \textbf{What the API assumes about it} \\
\midrule
\endhead
SoC & SM8850P, ``Snapdragon 8 Elite Gen 5'', codename \code{canoe}, SoC ID 661.
  Android 16, userdebug, \code{su 0} available. \\
CPU & Oryon-2, 6 Gold + 2 Prime. \textbf{Gold is cores 0--5} (28-point OPP
  ladder topping out at 3.6288\,GHz); \textbf{Prime is cores 6--7} (27 points to
  4.3968\,GHz). The two ladders share only four setpoints. These ceilings are
  hard-coded as \code{CLUSTER\_MAX\_GHZ} and used to flag under-clocked runs
  (\S\ref{sec:clock}). \\
GPU & Adreno 840, OpenCL, 18 power levels from 160\,MHz to 1.2\,GHz. GPU
  hardware perfcounters are \emph{not} readable on this build; the GPU reports
  \code{CL\_PROFILING} durations only. \\
DSP & Hexagon cDSP v81, 8 HVX contexts, 8\,MB VTCM, reached over FastRPC. \\
NPU & Hexagon HTP, driven through QAIRT/QNN 2.40. Shares the Hexagon subsystem
  with the DSP, which has consequences for parallelism (\S\ref{sec:parallel}). \\
\bottomrule
\end{longtable}

\noindent \textbf{The Prime ceiling is 4.3968\,GHz, not 4.608.} The sysfs node
\code{cpuinfo\_max\_freq} reports 4608000 and the governor will accept it, but
that setpoint is absent from \code{scaling\_available\_frequencies}, the
effective clock only reaches about 3.8--3.9\,GHz there, and under sustained load
the governor settles to 4.34. 4.3968 is the fastest setpoint that can actually
be selected and held, so it is the one the API compares against.

\subsection{Why it is called an API}

The name is accurate in the ordinary sense: it is a programmatic interface with
a small, stable vocabulary of objects (\code{Task}, \code{Result},
\code{Profiler}, \code{Device}, \code{Backend}) that hides a transport and a
parsing layer behind method calls. Callers write Python against those objects;
they never construct shell strings, never touch \code{adb}, and never write a
per-backend output parser. Swapping the engine a workload runs on is a
one-token change, which is precisely the property an interface is supposed to
give you.

It is a \emph{host-side} API: it is imported into a Python process on your
workstation, and the boundary it wraps is the USB link to the board, not a
network service. There is no server and no HTTP.

\subsection{Design commitments}

\begin{description}[leftmargin=1.6cm,style=nextline,font=\normalfont\itshape]
  \item[Uniform, not lossy.] Every backend reports through the same
    \code{Result} shape, but nothing is averaged away or hidden. Every
    \code{key=value} token the kernel printed survives in \code{Result.fields}
    and the untouched stdout survives in \code{Result.raw}.
  \item[Tool, not agent.] The API does exactly what it is told. It never picks
    a variant for you, never retries, never warms up behind your back, never
    tunes, and never flags an anomaly. Defaults exist so a call is short; they
    are documented here, printed by \code{-{}-catalogue}, and visible in the
    command it builds (\code{verbose=True}).
  \item[Costs stay visible.] Backend-native compute time and host-side
    transport are reported as separate fields. The NPU's graph-compile time,
    the GPU's host-to-device copy and the DSP's FastRPC marshalling are not
    folded into the headline number.
  \item[Real concurrency.] Parallel execution is genuine simultaneous hardware
    use, not interleaving. See \S\ref{sec:parallel}.
  \item[Fail early, on the host.] A mistyped variant, an impossible size, or
    \code{threads=} on a kernel that cannot take it raises \code{ValueError} at
    \code{Task} construction --- not on the device three dispatches into a
    sweep. The messages name the valid alternatives.
  \item[Refuse rather than guess.] Where a number cannot be produced honestly
    the API says so. \code{power()} returns \code{available: False} with a
    reason for the DSP rather than inventing watts; a counter this silicon does
    not implement raises instead of returning the silent zero the hardware
    would give you.
  \item[Standalone.] The module depends only on the Python standard library and
    an \code{adb} on \code{PATH}. It does not require the control container,
    the storage stack, or the orchestrator.
\end{description}

\subsection{What it deliberately does not do}

It does not build or push kernels (see \S\ref{sec:install}), does not persist
results to the project's Parquet/SQLite store, does not capture Perfetto
traces, does not read GPU, DSP or NPU hardware counters (only CPU PMU counters,
\S\ref{sec:counters}), and does not measure energy directly --- there is no
readable power rail on this board, so \code{power()} is a fitted model and says
so (\S\ref{sec:power}). It is the dispatch and measurement front door; the
heavier subsystems live elsewhere in the repository.

% ============================================================================
\section{Installation and prerequisites}
\label{sec:install}
% ============================================================================

\subsection{Requirements}

\begin{longtable}{@{}p{4.2cm}p{10.5cm}@{}}
\toprule
\textbf{Requirement} & \textbf{Notes} \\
\midrule
\endhead
Python 3.10 or newer & The module uses PEP 604 unions (\code{float | None})
  and builtin generics. 3.9 will fail to import. No \code{pip install} is
  needed: the module imports only from the standard library. \\
\code{adb} on \code{PATH} & Any recent platform-tools. Verify with
  \code{adb devices}; exactly one device should be listed, or pass a serial
  explicitly. Version 30 or newer is recommended --- an \code{adb} older than
  the device's \code{adbd} negotiates badly. \\
An HDK8850 (SM8850P) board & Connected over USB-C, in \code{device} state, with
  USB debugging authorised. TCP transport works too if \code{adb connect} has
  already been done. \\
Kernels present on device & The prebuilt binaries must exist under
  \code{/data/local/tmp/profiler-kernels}. \\
\code{su 0} on the device & Required \emph{only} for hardware counters
  (\S\ref{sec:counters}). Everything else runs at shell uid. \\
\bottomrule
\end{longtable}

\subsection{Provisioning the device}

The wrapper assumes the kernel tree already sits on the board. Push it with the
project's provisioner, which also verifies five functional gates before it
reports success:

\begin{lstlisting}[style=sh]
python scripts/provision_device.py \
    --qairt-root "<path>/qairt/2.40.0.251030" \
    --dsp-skel   kernels/matmul/dsp/prebuilt/libdsp_matmul_skel.so
python scripts/verify_board.py        # expect: 0 fail -> BOARD OK.
\end{lstlisting}

\noindent To go the other way --- put the board back to unprovisioned, so
the setup path can be tested from scratch --- use \code{scripts/reset\_device.py}.
It reports what it would delete and requires \code{-{}-yes} before
removing anything. \textbf{Rebooting does not do this}: \code{/data} is f2fs
on a real block device, so \code{/data/local/tmp} survives a reboot and
the kernels are still there afterwards.

\noindent Provisioning is idempotent: re-running after a failure is the
normal remedy.

\noindent The expected on-device layout is:

\begin{lstlisting}[style=plain]
/data/local/tmp/profiler-kernels/
    matmul_neon_handtuned, matmul_eigen_baseline   <- CPU matmul at the base
    gpu/         matmul_opencl_*                   + libs (LD_LIBRARY_PATH=".")
    dsp/         *_dsp*, stubs and skels           (FastRPC)
    npu/         *_qnn_*                           + QNN backend libraries
    saxpy/       saxpy_scalar, saxpy_neon, saxpy_opencl_fp32
    reduction/   reduction_scalar, reduction_neon, reduction_opencl_fp32
    convolution/ conv_scalar, conv_neon, conv_int8_neon, conv_opencl_fp32
    nbody/       nbody_neon, nbody_opencl_fp32
    decode/      decode_neon, decode_opencl_fp32, decode_opencl_gemv
    cachelat/    cache_latency
    memlab/      gpu_memlab
    custom/      whatever you build yourself
/data/local/tmp/
    mnv2/  nlp/  qwen/  llama/  sd15/               <- converted networks (DLCs)
\end{lstlisting}

\noindent If a binary is missing, the corresponding call fails with a shell
error in \code{Result.error} rather than raising. Networks are staged
separately and are not part of provisioning; see \S\ref{sec:models}.

\subsection{Smoke test}

\begin{lstlisting}[style=sh]
python profiler.py --demo
\end{lstlisting}

\noindent This runs a short live tour: a connectivity check, one convolution on
all four backends in parallel, an NPU roofline inversion pair, a CPU affinity
comparison, a hardware-counter comparison at one and six threads, and a small
dependency DAG. If the demo prints results, the API is correctly wired to your
board.

% ============================================================================
\section{Quick start}
% ============================================================================

\begin{lstlisting}
from profiler import Profiler, Backend, matmul, saxpy, reduction, convolution

p = Profiler()                                   # uses the only attached device

# 1. One dispatch.
r = p.run(convolution(Backend.NPU, cin=8, cout=8))
print(r.summary())
print(r.compute_us, r.gflops, r.fields["finalize_us"])

# 2. The same workload on every engine, at the same time.
for r in p.run_parallel([convolution(b, cin=8, cout=8) for b in Backend]):
    print(r.summary())

# 3. The same batch, strictly one at a time (clean, uncontended numbers).
for r in p.run_serial([convolution(b, cin=8, cout=8) for b in Backend]):
    print(r.summary())

# 4. Hardware counters riding along with a CPU dispatch.
r = p.run(convolution(Backend.CPU, cin=32, cout=32,
                      threads=6, core="gold",          # where, and how many
                      counters="cache", per_core=True))
print(r.derived["ipc"], r.derived["l1d_miss_pct"])
print(r.counter_table())

# 5. Dependency order: warm up on CPU, then fan out to GPU and NPU.
results = p.run_dag([
    reduction(Backend.CPU, size=1 << 20, name="warmup"),
    convolution(Backend.GPU, cin=16, cout=16, name="gpu", depends_on=["warmup"]),
    convolution(Backend.NPU, cin=16, cout=16, name="npu", depends_on=["warmup"]),
])
print(results["npu"].summary())
\end{lstlisting}

\noindent Everything else in this document is detail on those five calls, plus
two subsystems that stand slightly apart: converted networks
(\S\ref{sec:models}) and power estimation (\S\ref{sec:power}).

% ============================================================================
\section{Execution model}
% ============================================================================

\subsection{The path of one dispatch}

\begin{lstlisting}[style=plain]
  Task(workload, backend, params, variant, core, threads, counters, ...)
        |
        |  1. registry lookup: WORKLOADS[workload].kernels[backend] -> KernelEntry
        |  2. variant -> binary name (or a --variant flag, DSP and nbody-GPU)
        |  3. entry.defaults overlaid with task.params       (task wins)
        |  4. entry.param_map renames keys the kernel spells differently
        |  5. params -> CLI flags   (perf_mode="sustained" -> --perf-mode sustained)
        |  6. backend prologue: cd / LD_LIBRARY_PATH / ADSP_LIBRARY_PATH / taskset
        |  7. if counters: prepend `simpleperf stat --csv -e ...`
        |                  and wrap the WHOLE command in `su 0 sh -c '...'`
        v
  build_command(task, base)  ->  one shell string
        |
        v
  Device.shell(cmd, timeout)  ->  adb [-s SERIAL] shell "<cmd>"   (blocking)
        |
        v
  stdout scraped twice:  key=value / "key: value"  -> Result.fields
                         simpleperf --csv rows      -> Result.counters/.per_core
        |
        v
  Result(ok, compute_us, wall_us, gflops, gbps, fingerprint, fields,
         counters, derived, per_core, raw)
\end{lstlisting}

\noindent One \code{Task} is exactly one process on the device. Nothing is
batched, cached or reordered. The two parsers read the same stream without
colliding, because simpleperf's CSV rows are comma-separated and the kernels'
output is not.

\subsection{Backend command construction}

Each engine needs a different environment. The API applies it for you:

\begin{longtable}{@{}p{1.6cm}p{13.1cm}@{}}
\toprule
\textbf{Backend} & \textbf{Command shape} \\
\midrule
\endhead
CPU & \code{taskset <mask> <base>/<subdir>/<binary> <flags>} \newline
  The affinity mask comes from \code{Task.core} (\S\ref{sec:affinity}). \\
GPU & \code{cd <dir> \&\& LD\_LIBRARY\_PATH=<ldpath> ./<binary> <flags>} \newline
  Adreno matmul and decode ship their libraries beside the binary
  (\code{ldpath="."}); the other GPU kernels use \code{/vendor/lib64}. \\
DSP & \code{cd <dir> \&\& LD\_LIBRARY\_PATH=. ADSP\_LIBRARY\_PATH='.;<dir>' ./<binary> <flags>} \newline
  FastRPC needs the stub next to the working directory and the skel on the
  DSP's library path. \\
NPU & \code{cd <dir> \&\& LD\_LIBRARY\_PATH=<dir> ADSP\_LIBRARY\_PATH=<dir> ./<binary> <flags>} \newline
  QNN resolves its HTP backend libraries from those paths. \\
model & \code{cd <modeldir> \&\& rm -rf <out> \&\& LD\_LIBRARY\_PATH=<modeldir> \dots{}
  ./qnn-net-run -{}-dlc\_path \dots{} -{}-backend libQnn\{Cpu,Gpu,Htp\}.so \dots} \newline
  Networks bypass the kernel tree entirely (\S\ref{sec:models}). \\
\bottomrule
\end{longtable}

\noindent Pass \code{verbose=True} to the \code{Profiler} to print each command
before it runs --- the fastest way to confirm what the API actually did --- or
call \code{build\_command} to get the string without a board attached.

\subsection{Privilege: why counters run under \texttt{su 0}}
\label{sec:su}

When, and only when, counters are requested, the entire command is re-wrapped:

\begin{lstlisting}[style=sh]
su 0 sh -c '<the command that would otherwise have run>'
\end{lstlisting}

On this board the shell domain cannot open the PMU. \code{simpleperf} answers
\emph{``Event type `cpu-cycles' is not supported on the device''} for every
event, including ones \code{simpleperf list} advertises, while the identical
command under \code{su 0} counts normally. \code{perf\_event\_paranoid} is
already \code{-1}, so this is an SELinux domain restriction, not a paranoia
level.

\textbf{This is not \code{adb root}, and the difference matters.} \code{su 0}
elevates one command from an unrooted \code{adbd}. \code{adb root} restarts
\code{adbd} itself in a domain that cannot reach the OpenCL driver, which
breaks every GPU dispatch --- \code{clGetPlatformIDs} returns \code{-1001}, no
platform. Never run \code{adb root} on this board while profiling; the API is
careful not to need it.

\subsection{Serial, parallel, and DAG semantics}
\label{sec:parallel}

\begin{description}[leftmargin=1.8cm,style=nextline,font=\normalfont\itshape]
  \item[Serial] Tasks run strictly one after another. Nothing else the API
    controls is on the machine while a measurement is taken. Use this for any
    number you intend to quote.
  \item[Parallel] Every task is launched on its own host thread, so $N$ blocking
    \code{adb shell} invocations are in flight at once. The CPU cluster, the
    Adreno GPU and the Hexagon subsystem are independent hardware, so tasks on
    those genuinely execute simultaneously. Tasks on the \emph{same} backend
    contend --- that contention is real and measurable, and is often the point,
    but it is not a clean single-engine number.

    \textbf{One exception, and it is not visible from the host: NPU and DSP do
    not run in parallel.} They are two software paths --- QNN HTP and FastRPC ---
    onto the \emph{same} Hexagon cDSP, so they serialise in hardware. Submitting
    both concurrently looks correct from here: the dispatches go out together and
    their host-side windows overlap. The cDSP runs them one after the other
    anyway. Measured in Discovery~50, which found NPU+DSP to be the sole
    \textsc{hw-serialize} pair in the nine-pair matrix; first caught when a
    \code{wait\_any([dsp, npu])} returned the DSP result while the NPU sat
    blocked behind cDSP warm-up. Effective width is four independent paths, or
    three whenever NPU and DSP are both in play. Do not read an NPU+DSP parallel
    batch as a speed-up measurement.

    \textbf{The same sharing cuts the other way.} A busy NPU raises the Hexagon
    DVFS operating point, and the DSP --- which cannot vote for its own clock on
    this board (\S\ref{sec:dspclock}) --- is carried along with it. A DSP kernel
    measured next to an NPU peer can be \emph{faster} than the same kernel
    measured alone (Discovery~75: 1.59$\times$). Every DSP figure in this
    project is therefore the idle-Hexagon case, and ``isolation'' is not the
    cleanest measurement on this engine.
  \item[DAG] Tasks are grouped into levels: every task whose dependencies have
    completed runs in the same parallel batch, then the next level starts. A
    task never starts before its dependencies finish. Cycles and unknown
    dependency names raise \code{ValueError} before anything is dispatched.
\end{description}

\noindent Wall-clock overlap in parallel mode is host-scheduled, so start times
are not synchronised to the microsecond. For tightly-coupled interference
studies, use longer per-kernel iteration counts so the overlap window dominates
the launch skew.

% ============================================================================
\section{API reference}
% ============================================================================

% ----------------------------------------------------------------------------
\apient{Backend}

An \code{Enum} of the four compute engines. \code{str(Backend.CPU) == "cpu"},
and the enum is iterable, so \code{for b in Backend} sweeps all four.

\begin{longtable}{@{}llp{9.4cm}@{}}
\toprule
\textbf{Member} & \textbf{Value} & \textbf{Hardware and native timing source} \\
\midrule
\endhead
\code{Backend.CPU} & \code{"cpu"} & Oryon-2 cores, affinity via \code{taskset};
  self-timed \code{wall\_us} or \code{mean\_us}. \\
\code{Backend.GPU} & \code{"gpu"} & Adreno 840, OpenCL; \code{CL\_PROFILING}
  \code{duration\_ns}. \\
\code{Backend.DSP} & \code{"dsp"} & Hexagon cDSP v81, FastRPC + HVX;
  on-DSP \code{dsp\_us}. \\
\code{Backend.NPU} & \code{"npu"} & Hexagon HTP via QNN; \code{execute\_us\_mean}
  plus a separate \code{finalize\_us}. \\
\bottomrule
\end{longtable}

\noindent Not every workload exists on every backend. \code{nbody} is CPU and
GPU only; \code{cachelat} is CPU only; \code{gpu\_memlab} is GPU only; \code{decode} and \code{model} skip the
DSP. Asking for a pair that does not exist raises \code{ValueError} naming the
backends that do.

% ----------------------------------------------------------------------------
\apient{Task}
\label{sec:task}

A description of one dispatch. Constructing a \code{Task} performs no I/O; it
validates the workload/backend pair, the variant, the size, the thread request
and the counter selection, then assigns a name.

\begin{lstlisting}
Task(workload: str,
     backend: Backend,
     params: dict = {},
     variant: str | None = None,
     name: str | None = None,
     timeout: float = 180.0,
     depends_on: list[str] = [],
     core: object = None,
     threads: int | None = None,
     counters: object = None,
     per_core: bool = False,
     counter_cpus: object = None,
     binary: str | None = None)
\end{lstlisting}

\begin{longtable}{@{}p{2.6cm}p{2.5cm}p{9.6cm}@{}}
\toprule
\textbf{Field} & \textbf{Type} & \textbf{Meaning} \\
\midrule
\endhead
\code{workload} & \code{str} & A key of \code{WORKLOADS}: \code{"matmul"},
  \code{"saxpy"}, \code{"reduction"}, \code{"convolution"}, \code{"nbody"},
  \code{"decode"}, \code{"cachelat"}, \code{"morph"}, \code{"gpu\_memlab"}, \code{"custom"}, \code{"model"}. Unknown values raise
  \code{ValueError} listing the known ones. \\
\code{backend} & \code{Backend} & Which engine. A pair with no kernel raises
  \code{ValueError} naming the backends that workload does have. \\
\code{params} & \code{dict} & Kernel parameters, passed through as CLI flags
  (\S\ref{sec:params}). \\
\code{variant} & \code{str | None} & Which kernel implementation
  (\S\ref{sec:catalogue}). \code{None} selects the entry's default. A name that
  does not exist raises \code{ValueError} listing every valid variant, the
  default, and --- when the typo is close to a real name --- a
  ``did you mean'' line. \\
\code{name} & \code{str | None} & Identifier used by \code{run\_dag} and printed
  by \code{Result.summary()}. Defaults to
  \code{"<workload>.<backend>[.<variant>]"}. Must be unique within a DAG. \\
\code{timeout} & \code{float} & Seconds before the \code{adb} call is abandoned
  and the task fails with exit code 124. Raise it for NPU work
  (\S\ref{sec:caveats}). \\
\code{depends\_on} & \code{list[str]} & Names of tasks that must complete first.
  Only consulted by \code{run\_dag}. \\
\code{core} & see \S\ref{sec:affinity} & CPU affinity: \emph{where} the work is
  allowed to run. Ignored on non-CPU backends. \\
\code{threads} & \code{int | None} & \emph{How many} worker threads the kernel
  spawns, passed as \code{-{}-threads N}. Only accepted by threaded variants
  (\S\ref{sec:threads}). \\
\code{counters} & see \S\ref{sec:counters} & CPU hardware counters to collect
  alongside the dispatch. \code{None} means none. \\
\code{per\_core} & \code{bool} & Break the counter totals down per CPU
  (\code{simpleperf -{}-per-core}). Only meaningful with \code{counters}. \\
\code{counter\_cpus} & \code{int | list | str | None} & Count \emph{per-CPU and
  system-wide} on these CPUs instead of per-process. See
  \S\ref{sec:countercpus}. \\
\code{binary} & \code{str | None} & Dispatch this device binary instead of
  looking one up in the registry. Used by \code{custom()}
  (\S\ref{sec:custom}). \\
\bottomrule
\end{longtable}

\subsubsection*{Read-only properties}

\begin{longtable}{@{}p{3.4cm}p{11.4cm}@{}}
\toprule
\textbf{Property} & \textbf{Value} \\
\midrule
\endhead
\code{task.entry} & The \code{KernelEntry} this task resolves to: device
  subdirectory, variant table, always-on flags. \\
\code{task.resolved\_variant} & The variant that will actually run ---
  \code{variant} if given, else the registry default. Worth checking: defaults
  are deliberately \emph{not} the fastest variant, so that the naive/optimised
  gap stays visible. \\
\code{task.counter\_events} & The flat, deduplicated event list that
  \code{counters} resolves to. Recomputed on every access, so assigning
  \code{task.counters = "stalls"} after construction takes effect. \\
\code{task.cores} & The list of CPU ids the affinity mask actually permits,
  e.g.\ \code{[6, 7]} for \code{core="prime"}. \\
\code{task.oversubscribed} & \code{True} when \code{threads} exceeds
  \code{len(task.cores)} --- more workers than cores they are allowed on. This
  is legal and measurable, not an error. \\
\bottomrule
\end{longtable}

\subsubsection*{What is validated at construction}

All five checks run before any I/O, so a sweep fails on line one rather than
twenty minutes in:

\begin{enumerate}[leftmargin=0.8cm]
  \item the workload exists;
  \item the workload has a kernel on that backend;
  \item the variant exists on that entry;
  \item \code{size} matches \code{fixed\_size} when the chosen variant is
    compiled for exactly one shape (\S\ref{sec:fixedsize});
  \item \code{threads} is $\geq 1$ \emph{and} the resolved variant is in
    \code{threaded\_variants};
  \item every requested counter is implemented on this silicon.
\end{enumerate}

\subsubsection*{Convenience constructors}

Eight module-level functions build a \code{Task} without naming the workload
string. Task-level keywords (\code{name}, \code{variant}, \code{depends\_on},
\code{timeout}, \code{core}, \code{threads}, \code{counters}, \code{per\_core},
\code{counter\_cpus}) are recognised and peeled off; \emph{every other keyword
becomes a kernel parameter}.

\begin{lstlisting}
matmul(backend, **params)          -> Task   # dense square GEMM
saxpy(backend, **params)           -> Task   # y = a*x + y
reduction(backend, **params)       -> Task   # sum of floats
convolution(backend, **params)     -> Task   # 2D conv
nbody(backend, **params)           -> Task   # direct-summation N-body (CPU/GPU)
decode(backend, **params)          -> Task   # weight-stationary GEMM
custom(backend, binary, **params)  -> Task   # a kernel you wrote
model(backend, name, **params)     -> Task   # a converted network

# Equivalent:
saxpy(Backend.CPU, size=1 << 20, variant="scalar", core="prime", name="s1")
Task("saxpy", Backend.CPU, {"size": 1 << 20}, "scalar", "s1", core="prime")
\end{lstlisting}

\begin{longtable}{@{}p{2.3cm}p{12.5cm}@{}}
\toprule
\textbf{Constructor} & \textbf{What the workload is for} \\
\midrule
\endhead
\code{matmul} & Dense square matmul, \code{size=N} (default 512). Arithmetic
  intensity $\approx N/6$: the compute-bound end of the set. Variants differ
  enormously --- GPU \code{best} is about 16$\times$ \code{naive}, CPU
  \code{eigen} about 3$\times$ \code{neon}. \\
\code{saxpy} & $y = ax + y$ over \code{size} floats (default 1M). AI 0.17 ---
  the memory-bound end, and the workload where one CPU cluster beats the NPU by
  about 50$\times$. \\
\code{reduction} & Sum of \code{size} floats (default 4M). Memory-bound like
  saxpy, \emph{but} the NPU has a native \code{ReduceSum} so it compiles in
  milliseconds rather than seconds. The saxpy/reduction pair is a controlled
  experiment in operator support, holding arithmetic intensity roughly fixed. \\
\code{convolution} & 2D convolution over \code{cin}/\code{cout}/\code{H}/\code{W}/\code{k}.
  AI $\approx 18$ at $C=8$, enough to feed the NPU's tensor cores: it wins by
  about 68$\times$ over one CPU cluster. \\
\code{nbody} & Direct-summation N-body, $O(N^2)$ arithmetic on $O(N)$ data, so
  AI $= 1.25N$ and \emph{grows with the problem}. QNN cannot express it, so
  this is the one workload where the GPU competes on equal terms. CPU and GPU
  only, by design rather than omission. \\
\code{decode} & Weight-stationary GEMM $C[S \times n] = A[S \times k] \cdot
  W[k \times n]$ with the token count \code{tokens} exposed. That one knob
  walks the workload across the roofline ridge: $S=1$ is LLM decode (every
  weight read once, DRAM-bound), large $S$ is prefill (weights reused,
  compute-bound). The ranking inverts across it --- at $S=1$ one CPU core beats
  the NPU; at $S=256$ the NPU wins 29$\times$ (Discovery~66). \\
\code{custom} & Your own kernel, by device filename. \S\ref{sec:custom}. \\
\code{model} & A converted network through \code{qnn-net-run}.
  \S\ref{sec:models}. \\
\bottomrule
\end{longtable}

\noindent \textbf{Unrecognised keywords silently become flags.} A typo such as
\code{sizee=1024} produces \code{-{}-sizee 1024} on the device rather than a
Python error. Most kernels reject unknown flags and the task then fails, so
check \code{Result.error} first when a call fails unexpectedly --- but some
kernels ignore them, in which case you have quietly measured the kernel's
\emph{default} instead of what you asked for. There is no way for the host to
tell those two apart, which is why \code{verbose=True} is worth having on while
you develop a sweep.

% ----------------------------------------------------------------------------
\apient{CPU affinity: the \texttt{core} argument}
\label{sec:affinity}

\code{core} decides \emph{where} work may run. It accepts four forms and is
resolved to a \code{taskset} hexadecimal mask.

\begin{longtable}{@{}p{3.6cm}p{1.6cm}p{9.5cm}@{}}
\toprule
\textbf{Form} & \textbf{Mask} & \textbf{Meaning} \\
\midrule
\endhead
\code{None} (default) & \code{f0} & Cores 4--7. This is what the project's sweep
  scripts use, so defaults reproduce the published roofline numbers.
  \textbf{Read the note below before relying on it.} \\
\code{core=7} & \code{80} & That one core. \\
\code{core=[6, 7]} & \code{c0} & Those cores, OR-ed together. \\
\code{core="prime"} & \code{c0} & Prime cluster, cores 6--7, up to 4.3968\,GHz. \\
\code{core="gold"} & \code{3f} & Gold cluster, cores 0--5, up to 3.6288\,GHz. \\
\code{core="all"} & \code{ff} & All eight cores. \\
\code{core="0xc0"} / \code{"c0"} & as given & Raw mask, passed through
  unchanged and not validated. \\
\bottomrule
\end{longtable}

\noindent\textbf{The default mask straddles the cluster boundary.} On this SoC
Gold is cores 0--5 and Prime is 6--7, so \code{f0} (cores 4, 5, 6, 7) permits
\emph{two Gold cores and two Prime cores}. The scheduler may place a
single-threaded kernel on either, and two runs with identical arguments can
land on different clusters with a 1.13$\times$ clock difference between them.
Two consequences:

\begin{itemize}[leftmargin=0.6cm]
  \item For a number you intend to quote, name the cluster explicitly:
    \code{core="prime"} or \code{core="gold"}.
  \item \code{clock\_of\_max\_pct} and \code{below\_max\_clock}
    (\S\ref{sec:clock}) are computed \emph{only} when the mask lies entirely
    inside one cluster, because only then is there an unambiguous ceiling to
    compare against. Under the default mask they are silently absent.
\end{itemize}

\noindent Cluster membership is fixed for the SM8850P; masks are not
re-derived at runtime.

% ----------------------------------------------------------------------------
\apient{Threading: the \texttt{threads} argument}
\label{sec:threads}

\code{threads=N} passes \code{-{}-threads N} to the kernel, asking it to spawn
$N$ workers. It is independent of \code{core}, which only decides which CPUs
those workers are \emph{allowed} on. Both are recorded, and the combination is
free: six threads on a one-core mask is a legal, measurable experiment ---
\code{task.oversubscribed} reports \code{True} and the kernel time-slices.

\textbf{Threading is per-variant, not per-workload.} The binary either accepts
the flag or hard-errors on it, and variants inside one registry entry differ:
\code{conv\_neon} is threaded, \code{conv\_scalar} sits in the same entry and
rejects \code{-{}-threads}. Asking for threads on a variant that cannot take
them raises at construction and names the ones that can.

\begin{longtable}{@{}p{3.2cm}p{11.6cm}@{}}
\toprule
\textbf{Workload} & \textbf{Threaded variants (CPU)} \\
\midrule
\endhead
\code{matmul} & \code{neon}. \emph{Not} \code{eigen}. \\
\code{saxpy} & \code{scalar}, \code{neon} \\
\code{reduction} & \code{scalar}, \code{neon} \\
\code{convolution} & \code{neon}, \code{int8}. \emph{Not} \code{scalar}. \\
\code{nbody} & \code{neon} \\
\bottomrule
\end{longtable}

\noindent \textbf{A trap worth knowing about matmul.} The CPU matmul spawns and
joins its threads \emph{inside} the timed region, so a single-iteration run
charges the whole thread-spawn cost to a short measurement --- and the shorter
the region, the worse parallelism looks. Measured on Gold, matmul $N=512$
scales 2.55$\times$ at one iteration and 4.11$\times$ at thirty. The registry
therefore sends \code{-{}-iterations 10} by default for CPU matmul. If you are
studying parallel efficiency, raise it further; if you lower it to 1 you will
measure thread creation and call it poor scaling.

% ----------------------------------------------------------------------------
\apient{Hardware counters: \texttt{counters}, \texttt{per\_core}, \texttt{counter\_cpus}}
\label{sec:counters}

Any CPU dispatch can carry ARMv9 PMU counters, collected by \code{simpleperf}
wrapped around the same command.

The prefix is applied on \emph{every} backend, not only the CPU, and a GPU,
DSP or NPU dispatch with counters succeeds. What comes back is the counters of
the 	extbf{ARM process driving the accelerator}, not of the accelerator ---
which is a useful measurement, and not the one people assume. See
\S
ef{sec:accelcounters}. There is no way to read Adreno or Hexagon hardware
counters from this board at all.

\begin{lstlisting}
r = p.run(convolution(Backend.CPU, cin=32, cout=32, core="gold", threads=6,
                      counters="cache", per_core=True))

r.counters    # {'cpu-cycles': 1.9e10, 'raw-l1d-cache': 1.2e8, ...}  raw counts
r.derived     # {'ipc': 2.91, 'l1d_miss_pct': 0.28, 'effective_ghz': 3.41, ...}
r.per_core    # {0: {...}, 1: {...}, ...}  one dict per CPU
print(r.counter_table())
\end{lstlisting}

\subsubsection*{\texttt{counters} --- what to collect}

Accepts \code{None}, a preset name, a single event name, or a list mixing the
two. Presets expand; duplicates collapse; order is preserved.

\begin{longtable}{@{}p{1.9cm}p{1.1cm}p{11.6cm}@{}}
\toprule
\textbf{Preset} & \textbf{Events} & \textbf{What it answers} \\
\midrule
\endhead
\code{basic} & 2 & Cycles and instructions --- IPC, and nothing else. The
  cheapest way to get \code{effective\_ghz}. \\
\code{branch} & 7 & Prediction accuracy, split by branch kind: direct, indirect
  (the expensive kind), and returns. \\
\code{cache} & 10 & L1D, L1I and L2 accesses, refills and writebacks. The
  default choice for a memory-hierarchy question. \\
\code{tlb} & 10 & L1/L2 TLB accesses and refills, plus page-table walks. \\
\code{stalls} & 8 & Frontend vs backend stalls, in both cycles and issue slots,
  plus memory-attributed backend stalls. \\
\code{mix} & 9 & Instruction mix: retired vs speculative, loads, stores,
  integer, and FP/SIMD. \\
\code{memory} & 9 & Memory and bus accesses with the refills and memory stalls
  that explain them. \\
\code{prefetch} & 9 & Cache traffic split by who generated it --- the program or
  the hardware prefetcher. \\
\bottomrule
\end{longtable}

\noindent \code{python profiler.py -{}-list-counters} prints all eight presets,
all 92 verified events grouped by subsystem, and all 43 blocked ones with the
reason for each.

\subsubsection*{The 12-counter limit}

The Oryon-2 PMU has \textbf{12 hardware counters per core}
(\code{MAX\_HW\_COUNTERS}). Up to twelve events are collected simultaneously and
exactly. Ask for more and \code{simpleperf} time-multiplexes them: each event is
counted for part of the run and scaled up, so the numbers become estimates.
The API does not stop you --- multiplexing is sometimes the right trade --- but
it sets \code{Result.multiplexed = True} and \code{summary()} prints
\code{[multiplexed]}. Every shipped preset is $\leq$ 12 events, so presets never
multiplex.

\subsubsection*{Blocked counters, and why refusing beats returning zero}

\code{simpleperf list raw} advertises 742 events on this part and \emph{none of
them error}. The ones the silicon does not implement return a silent, confident
\textbf{zero}. A zero is indistinguishable from a real measurement of nothing,
so it is the most dangerous possible failure mode: you can build an entire
argument on a counter that was never counting.

The module therefore ships a whitelist. 92 events were each run against a
memory-heavy workload on both clusters and returned non-zero; 43 more returned
exactly zero and are blocked. Requesting a blocked event raises
\code{ValueError} naming the reason:

\begin{lstlisting}[style=plain]
ValueError: counter 'raw-l3d-cache' reads a constant 0 on this board -- no
L3/system-cache PMU: the SLC is invisible to the core PMU on this SoC (sysfs
exposes only L1D/L1I/L2). It would not error on device, it would silently
return 0.
\end{lstlisting}

\noindent The four blocked families, and what to use instead:

\begin{longtable}{@{}p{4.3cm}p{10.4cm}@{}}
\toprule
\textbf{Blocked family} & \textbf{Reason, and the substitute} \\
\midrule
\endhead
\code{raw-l3d-*}, \code{raw-ll-*} & There is no L3/system-cache PMU: the SLC is
  invisible to the core PMU. L2 refills (\code{raw-l2d-cache-refill}) are the
  last visible level. \\
\code{raw-stall-backend-l1d}, \code{-l2d}, \code{-tlb}, and the frontend
  equivalents & Per-structure stall attribution is not implemented. Use
  \code{raw-stall-backend-mem} for memory-attributed stalls. \\
\code{raw-l1d-cache-hit}, \code{raw-l2d-cache-hit}, \code{raw-l1d-cache-allocate}
  & Cache \emph{hit} counters are not implemented. Derive hits as
  (accesses $-$ refills). \\
Prefetch refill/hit attribution (\code{*-refill-hwprf}, \code{*-hit-rd-f*}) &
  Only the prefetcher's \emph{access} counters work
  (\code{raw-l1d-cache-hwprf}, \code{raw-l2d-cache-hwprf}). You can count
  prefetcher traffic; you cannot measure its hit rate. \\
\code{raw-ase-inst-spec} & Reads 0 even on NEON code. SIMD counts under
  \code{raw-vfp-spec}. \\
\bottomrule
\end{longtable}

\subsubsection*{\texttt{per\_core} --- the per-CPU breakdown}

\code{per\_core=True} adds \code{-{}-per-core} to \code{simpleperf}, which
reports each CPU separately instead of one aggregate. \code{Result.per\_core}
becomes \code{\{cpu\_id: \{event: count\}\}}, and because \code{simpleperf}
prints no aggregate row in this mode, \code{Result.counters} is reconstructed by
summing the cores.

This is how you see load imbalance. Six threads with one straggler look
identical in the totals and obvious per core.

\subsubsection*{\texttt{counter\_cpus} --- process-scoped vs system-wide}
\label{sec:countercpus}

By default counting is \textbf{process-scoped}: only your workload is counted,
wherever the scheduler puts it. That is what you want almost always.

\code{counter\_cpus} switches to \code{simpleperf -{}-cpu}, which counts
\textbf{per-CPU and system-wide} on the named CPUs --- so it also catches
whatever else Android decides to run there. It accepts an int, a list of ints,
or a raw string such as \code{"0-5"}. Use it when the question is ``what is
happening on these cores'' rather than ``what is my process doing''; interpret
the result knowing that the answer includes other people's work.

\subsubsection*{What counters cost}

\code{simpleperf} itself costs nothing measurable. The governor sampler that
rides along with it (\S\ref{sec:clock}) is not free: it forks \code{cat} and
\code{sleep} every 2\,ms, which is cheap when cores are idle and expensive when
they are not --- measured at $+1.5\%$ on one Prime core with the machine
otherwise idle, and $+31\%$ on a six-thread convolution occupying the whole
Gold cluster. That perturbation would land squarely on the multi-core scaling
experiments and read as poor parallel efficiency, so \textbf{the sampler only
runs for single-threaded dispatches}. With \code{threads > 1} you get counters
without \code{scaling\_cur\_freq\_ghz}.

% ----------------------------------------------------------------------------
\apient{Effective clock and throttling detection}
\label{sec:clock}

Two different clock numbers can appear in \code{Result.derived}, and they
routinely disagree. Both are present only when counters were requested.

\begin{longtable}{@{}p{4.0cm}p{10.7cm}@{}}
\toprule
\textbf{Field} & \textbf{What it is} \\
\midrule
\endhead
\code{effective\_ghz} & \textbf{Measured.} Cycles actually retired, divided by
  \code{threads}, divided by \code{simpleperf}'s own \emph{Total test time}.
  This is the clock the silicon really ran at. \\
\code{scaling\_cur\_freq\_ghz} & \textbf{Requested.} The maximum value the
  governor asked for, sampled from sysfs every 2\,ms \emph{during} the run. \\
\code{clock\_of\_max\_pct} & \code{effective\_ghz} as a percentage of the
  cluster ceiling. \\
\code{scaling\_cur\_freq\_samples} & How many governor samples were taken;
  reported so you can judge the one above. \\
\bottomrule
\end{longtable}

\noindent \code{Result.below\_max\_clock} is \code{True} when
\code{effective\_ghz} falls below \code{CLOCK\_FLAG\_RATIO} (0.80) of the
cluster ceiling, and \code{summary()} then appends
\code{[clock NN\% of max]}.

\textbf{This flag does not mean thermal throttling.} Under the shipping
\code{walt} governor a low clock is ordinary DVFS. It means ``this did not run
at full speed''; whether that was intended is yours to judge. What it reliably
catches is a sweep that silently drops from 4.40 to 1.37\,GHz half way through,
which otherwise reads as a workload effect.

Three constraints on these fields, each the result of getting it wrong first:

\begin{itemize}[leftmargin=0.6cm]
  \item \textbf{The denominator must be \code{simpleperf}'s window, not the
    kernel's \code{wall\_us}.} The counts cover the whole process --- startup,
    input generation, teardown --- while \code{wall\_us} covers only the timed
    region. Dividing whole-process cycles by region-only time reported
    17\,GHz on a 4.4\,GHz part.
  \item \textbf{\code{scaling\_cur\_freq\_ghz} needs a long enough run.} It is
    reported only when the workload ran at least 15\,ms and produced at least
    five samples. A 7\,ms workload yields one or two in-run samples out of
    thirty --- the rest are idle before and after --- and the maximum then reads
    \emph{lower} than the effective clock, teaching the exact opposite of the
    truth. Gating on sample count does not catch this; gating on runtime does.
  \item \textbf{\code{clock\_of\_max\_pct} needs a single-cluster mask.} With
    the default \code{core=None} it is absent (\S\ref{sec:affinity}).
\end{itemize}

% ----------------------------------------------------------------------------
\apient{Profiler}

The dispatcher.

\begin{lstlisting}
Profiler(device: Device | None = None, *, verbose: bool = False,
         thermal: bool = True, thermal_limit_c: float = 90.0)
\end{lstlisting}

\noindent \code{device} defaults to \code{Device()}, which targets the only
attached \code{adb} device. \code{verbose=True} prints each generated shell
command to stdout before it runs. \code{thermal} and \code{thermal\_limit\_c}
control the temperature check described in \S\ref{sec:thermal}.

\subsubsection*{\texttt{run(tasks, *, parallel=False)}}

The general entry point.

\begin{longtable}{@{}p{5.2cm}p{9.7cm}@{}}
\toprule
\textbf{Call} & \textbf{Returns} \\
\midrule
\endhead
\code{run(task)} & a single \code{Result} \\
\code{run([t1, t2, ...])} & \code{list[Result]}, executed serially, input order
  preserved \\
\code{run([t1, t2, ...], parallel=True)} & \code{list[Result]}, executed
  concurrently, input order preserved \\
\bottomrule
\end{longtable}

\noindent The \code{parallel} flag is the whole difference between an isolated
measurement and a contention measurement, which makes A/B studies a one-token
edit:

\begin{lstlisting}
clean     = p.run(batch, parallel=False)
contended = p.run(batch, parallel=True)
\end{lstlisting}

\subsubsection*{\texttt{run\_serial(tasks) -> list[Result]}}

Explicit serial form. Equivalent to \code{run(tasks)} on a list.

\subsubsection*{\texttt{run\_parallel(tasks) -> list[Result]}}

Explicit parallel form. Launches one host thread per task
(\code{ThreadPoolExecutor} sized to the batch), so all dispatches are in flight
simultaneously. Returns when the last one completes; results stay in input
order. An empty list returns an empty list.

\subsubsection*{\texttt{run\_dag(tasks) -> dict[str, Result]}}

Executes named tasks in dependency order, one parallel level at a time, and
returns a mapping from \code{Task.name} to \code{Result}.

\begin{description}[leftmargin=0.8cm]
  \item[Raises] \code{ValueError} if names are not unique, if a
    \code{depends\_on} entry names a task that is not in the list, or if the
    dependency graph contains a cycle. All three checks happen before any
    dispatch.
\end{description}

\noindent A failed dependency does \emph{not} cancel its dependents: a task
whose dependency returned \code{ok=False} still runs. Check \code{ok} on the
upstream result if that matters to your experiment.

\subsubsection*{\texttt{time\_model(backend, name, reps=100) -> dict}}
\label{sec:timemodel}

Per-inference time for a converted network, with graph preparation removed.
See \S\ref{sec:models}.

\subsubsection*{\texttt{power(task, *, seconds=35, artifact=None) -> dict}}

Estimated total-board power and energy while the task runs in a loop. See
\S\ref{sec:power}.

% ----------------------------------------------------------------------------
\apient{Thermal reporting}
\label{sec:thermal}

This HDK runs hot, and has fallen over: a sustained NPU graph reached
98\,\textdegree C with mitigation engaged, and the board dropped off USB and
needed its power cut by hand (Discovery~73). After every dispatch the profiler
reads the die sensors, records them on the \code{Result}, and prints one line
to stderr if the hottest zone has reached the limit (once per hot spell, again
only if it climbs another 5\,\textdegree C).

\begin{lstlisting}[style=plain]
  TOO HOT: 92 C on cpu-0-5-1 (limit 90 C). Let the board cool for a minute
  before the next run -- measurements taken now are throttled, and this board
  has dropped off USB at 98 C.
\end{lstlisting}

\noindent \code{summary()} also appends \code{[HOT 92C]}, and a later
dispatch that has come back down prints \code{cooled to 71 C.} once.

\subsubsection*{It reports; it does not act}

The profiler will not sleep, throttle, retry, re-order or refuse. Whether a hot
board invalidates \emph{your} measurement depends on what you are measuring ---
a thermal study wants the hot board --- so that judgement stays with you. This
is the same commitment as everywhere else in the API: it tells you what
happened and changes nothing on your behalf.

\subsubsection*{Where 90\,\textdegree C comes from}

Measured on board 5cabf4fd, 2026-09-22:

\begin{longtable}{@{}p{6.0cm}p{8.7cm}@{}}
\toprule
\textbf{State} & \textbf{Hottest zone} \\
\midrule
\endhead
Idle & $\approx$30\,\textdegree C \\
One heavy 6-thread CPU dispatch & 62\,\textdegree C \\
Four back to back (the lab plateau) & 77--79\,\textdegree C \\
All-core interference (Task 7), just after the run & 88--92\,\textdegree C, falls within seconds \\
Runaway NPU (Discovery~73) & 98\,\textdegree C, board lost \\
\bottomrule
\end{longtable}

\noindent 90 sits in the gap: clear of the 78\,\textdegree C an ordinary CPU
sweep reaches and of the brief peaks after an all-core run, so it does not cry
wolf, and still well under the 98 where this board has actually been lost. Change it with \code{thermal\_limit\_c}.

\subsubsection*{Which sensors, and the 105\,\textdegree C trap}

\code{/sys/class/thermal} exposes 135 entries on this board and \textbf{not
all of them are sensors}. \code{cpu-hw-trip-0} and \code{cpu-hw-trip-1} read a
constant \code{105000}: they are trip-point \emph{settings}. Taking the
maximum across every zone therefore reports 105\,\textdegree C forever, on an
idle board --- a false claim this project published once before anyone checked
the node type.

The selector is therefore a \textbf{whitelist}, \code{\_THERMAL\_ZONE\_RE},
covering 40 zones: 16 CPU core sensors, 4 cluster-cache, 11 \code{gpuss-*},
4 \code{nsphvx-*} (Hexagon HVX, the DSP), 4 \code{nsphmx-*} (Hexagon HMX, the
NPU) and \code{ddr}. Everything else --- trip points, battery-current-limit
levels, the \code{Debug\_Board} battery dummy, camera, modem and PMIC package
sensors --- is ignored. \code{verify\_board.py} uses the same whitelist.

\subsubsection*{What it costs, and when to switch it off}

Scanning all 135 zones costs about a second per read --- two \code{cat}s per
zone over \code{adb}. Instead the zone indices are resolved \textbf{once}
(about 560\,ms, on the first dispatch) and subsequent reads are a single
\code{cat} of the 40 whitelisted files.

\begin{longtable}{@{}p{4.6cm}p{10.1cm}@{}}
\toprule
\textbf{Setting} & \textbf{Cost, measured} \\
\midrule
\endhead
\code{thermal=True} (default) & about \textbf{16\,ms} added per dispatch in
  steady state, plus roughly 560\,ms once on the first. \\
\code{thermal=False} & nothing. \code{temp\_c} is \code{None} and
  \code{temps} is empty. \\
\bottomrule
\end{longtable}

\noindent Worth switching off for a long sweep of very short dispatches, where
16\,ms can exceed the kernel time and dominates the run. Leave it on for
anything that loads the board for more than a moment --- which is when it
matters.

\subsubsection*{Reading it yourself}

\begin{lstlisting}
r = p.run(convolution(Backend.CPU, cin=64, cout=64, threads=6, core="gold"))
r.temp_c                                   # 78.7
r.too_hot                                  # False at the default limit
sorted(r.temps.items(), key=lambda kv: -kv[1])[:3]
# [('cpu-0-5-1', 78.7), ('cpu-0-5-0', 75.7), ('cpu-0-1-1', 75.7)]
\end{lstlisting}

% ----------------------------------------------------------------------------
\apient{Result}
\label{sec:result}

The parsed outcome of one dispatch.

\begin{longtable}{@{}p{2.9cm}p{2.4cm}p{9.2cm}@{}}
\toprule
\textbf{Field} & \textbf{Type} & \textbf{Meaning} \\
\midrule
\endhead
\code{task} & \code{Task} & The task that produced this result. \\
\code{ok} & \code{bool} & \code{False} on non-zero exit, on a kernel-reported
  status that is not one of \code{OK}/\code{SUCCESS}/\code{PASS}/\code{PASSED}/\code{DONE},
  on timeout, or when the kernel printed no timing field at all. \\
\code{compute\_us} & \code{float | None} & Backend-native compute time in
  microseconds (\S\ref{sec:fields}). This is the number to compare across
  engines. \\
\code{compute\_key} & \code{str | None} & \emph{Which} kernel field
  \code{compute\_us} came from. Worth checking when comparing engines: the
  semantics differ per field. \\
\code{wall\_us} & \code{float | None} & The kernel's own wall-clock figure where
  it prints one. On the DSP and NPU this includes host-side transport, so
  \code{wall\_us} $\gg$ \code{compute\_us} is expected and meaningful. \\
\code{gflops} & \code{float | None} & Derived: FLOP count divided by
  \code{compute\_us}, with the iteration corrections of
  \S\ref{sec:iterations} applied. \\
\code{gbps} & \code{float | None} & Derived: byte count divided by
  \code{compute\_us}. Present when the kernel prints \code{bytes=}, when the
  workload has a \code{bytes\_fn}, or when the kernel prints its own
  \code{useful\_gbps}. \\
\code{fingerprint} & \code{float | None} & The kernel's numerical checksum, used
  to confirm the computation was correct before trusting its timing. For the
  QNN matmul clients, which print \code{fingerprint\_sum} and
  \code{fingerprint\_frob} but never a combined line, the two halves are added. \\
\code{fields} & \code{dict} & Every \code{key=value} token and \code{key: value}
  line parsed from stdout, numerics coerced to \code{int}/\code{float}, plus
  \code{host\_wall\_us} added by the profiler. \\
\code{raw} & \code{str} & Unmodified stdout. \\
\code{error} & \code{str | None} & Failure text (first 300 characters of stderr,
  or stdout, or the exit code) when \code{ok} is \code{False}. \\
\code{counters} & \code{dict} & Event \arrow{} raw count. Empty when no counters
  were requested. \\
\code{per\_core} & \code{dict} & \code{\{cpu\_id: \{event: count\}\}}. Empty
  unless \code{per\_core=True}. \\
\code{derived} & \code{dict} & Ratios the raw counts imply --- IPC, miss rates,
  stall percentages, effective clock. Only the ones the present events
  support. \\
\code{multiplexed} & \code{bool} & More than 12 events were requested, so the
  counts are scaled estimates rather than exact. \\
\code{below\_max\_clock} & \code{bool} & Ran under 80\% of the cluster ceiling
  (\S\ref{sec:clock}). \\
\code{temp\_c} & \code{float | None} & Hottest die sensor immediately
  \emph{after} the dispatch, in \textdegree C. \code{None} when thermal
  reporting is off (\S\ref{sec:thermal}). \\
\code{temps} & \code{dict} & Every zone read, \code{\{zone: degrees\}}. A
  78\,\textdegree C core beside a 30\,\textdegree C GPU is a different
  situation from both at 78. \\
\code{too\_hot} & \code{bool} & \code{temp\_c} reached the limit. Reported,
  never acted on. \\
\bottomrule
\end{longtable}

\subsubsection*{Properties}

\begin{longtable}{@{}p{4.0cm}p{10.7cm}@{}}
\toprule
\textbf{Property} & \textbf{Value} \\
\midrule
\endhead
\code{result.iterations} & How many times the kernel repeated the work inside
  one dispatch --- read from \code{iterations}/\code{reps}/\code{inner-reps} in
  the output, but only when \code{compute\_key} is a \emph{total} field
  (\code{wall\_us}). Otherwise 1. \\
\code{result.compute\_us\_per\_iter} & \code{compute\_us / iterations}: the
  per-dispatch latency. \\
\bottomrule
\end{longtable}

\subsubsection*{\texttt{summary() -> str}}

A one-line, column-aligned digest: task name; then either \code{N us} or
\code{N us/iter xM} when the kernel repeated the work $M$ times; then GFLOP/s,
GB/s, any of \code{ns\_per\_access} / \code{finalize\_us} / \code{transfer\_ns}
/ \code{readback\_ns} that are present, the fingerprint, IPC when counters ran,
and the \code{[multiplexed]} / \code{[clock NN\% of max]} markers. On failure it
returns \code{"[FAIL] <name>: <error>"}.

It is a display helper --- read the typed fields, not this string, when you are
computing on the results.

\subsubsection*{\texttt{counter\_table() -> str}}

A multi-line block: every raw count, then the per-core matrix if there is one,
then every derived ratio. Returns \code{"(no counters collected)"} when there
are none.

% ----------------------------------------------------------------------------
\apient{Device}

The \code{adb} transport.

\begin{lstlisting}
Device(serial: str | None = None, adb: str = "adb",
       base: str = "/data/local/tmp/profiler-kernels")

device.shell(command: str, timeout: float = 180.0) -> tuple[str, str, int]
device.is_connected() -> bool
\end{lstlisting}

\begin{itemize}[leftmargin=0.6cm]
  \item \code{serial} selects one board when several are attached
    (\code{adb devices}); \code{None} uses the only one.
  \item \code{adb} lets you point at a specific platform-tools binary.
  \item \code{base} relocates the on-device kernel tree. It does \emph{not}
    relocate converted networks, which live under \code{/data/local/tmp}
    independently.
  \item \code{shell} returns \code{(stdout, stderr, returncode)} and never
    raises: a timeout is reported as return code 124 with an explanatory
    stderr.
  \item \code{is\_connected} runs \code{echo ok} with a 10-second timeout.
\end{itemize}

\begin{lstlisting}
p = Profiler(Device(serial="5cabf4fd"), verbose=True)
if not p.device.is_connected():
    raise SystemExit("board not reachable")
\end{lstlisting}

% ----------------------------------------------------------------------------
\apient{build\_command(task, base) -> str}

Returns the exact device shell command a task would run, without running it.
Useful for logging a sweep's provenance, for pasting a single case into an
\code{adb shell} session, and for unit-testing flag construction with no
hardware attached.

\begin{lstlisting}
>>> from profiler import build_command, matmul, Backend, DEVICE_BASE
>>> print(build_command(matmul(Backend.DSP, size=512), DEVICE_BASE))
cd /data/local/tmp/profiler-kernels/dsp && LD_LIBRARY_PATH=.
ADSP_LIBRARY_PATH='.;/data/local/tmp/profiler-kernels/dsp'
./matmul_dsp_hvx --size 512

>>> print(build_command(matmul(Backend.GPU, variant="best"), DEVICE_BASE))
cd /data/local/tmp/profiler-kernels/gpu && LD_LIBRARY_PATH=.
./matmul_opencl_fp32_3x16u4_sg_half
\end{lstlisting}

\noindent (Line breaks above are for the page; the real return value is one
line.) With counters requested the returned string additionally carries the
\code{simpleperf} prefix, the governor sampler, and the surrounding
\code{su 0 sh -c '...'} wrapper.

% ----------------------------------------------------------------------------
\apient{The registry: \texttt{WORKLOADS}, \texttt{Workload}, \texttt{KernelEntry}}
\label{sec:registry}

\code{WORKLOADS} is a module-level dictionary mapping workload name to a
\code{Workload}, and is the single source of truth for the on-device layout. A
\code{Workload} carries its per-backend \code{KernelEntry} table plus optional
\code{flop\_fn} / \code{bytes\_fn} callables used to derive throughput for
kernels that do not print their own operation counts.

\begin{longtable}{@{}p{4.3cm}p{10.4cm}@{}}
\toprule
\textbf{\code{KernelEntry} field} & \textbf{Meaning} \\
\midrule
\endhead
\code{subdir} & Directory under the device base; \code{""} means the base
  itself. \\
\code{variants} & Variant name \arrow{} binary filename. \\
\code{default} & Variant used when the task does not name one. \\
\code{ldpath} & \code{LD\_LIBRARY\_PATH} for GPU dispatches. Default
  \code{/vendor/lib64}; \code{"."} where the kernel ships its own libraries. \\
\code{variant\_is\_flag} & \code{True} when the variant is selected by a
  \code{-{}-variant} flag on a single binary rather than by choosing a
  different file. True for the DSP saxpy/reduction/convolution entries and for
  GPU \code{nbody}. \\
\code{defaults} & Flags always applied unless the task overrides them, e.g.\
  \code{power=lab} on DSP, \code{perf\_mode=sustained} on NPU. \\
\code{fixed\_size} & Variant \arrow{} the only \code{-{}-size} it accepts
  (\S\ref{sec:fixedsize}). \\
\code{param\_map} & Rename a parameter on the way to the CLI, where the kernel
  spells it differently (\S\ref{sec:parammap}). \\
\code{threaded\_variants} & Which variants accept \code{-{}-threads}
  (\S\ref{sec:threads}). Empty means the whole entry is single-threaded. \\
\code{totals\_with\_per\_iteration\_time} & This kernel prints
  \code{flop=}/\code{bytes=} covering \emph{all} iterations while its timing
  field measures \emph{one} (\S\ref{sec:iterations}). \\
\bottomrule
\end{longtable}

\noindent Reading \code{WORKLOADS} is the supported way to enumerate what
exists:

\begin{lstlisting}
for name, wl in WORKLOADS.items():
    for backend, entry in wl.kernels.items():
        print(name, backend, sorted(entry.variants), "default:", entry.default)
\end{lstlisting}

% ============================================================================
\section{Kernel and variant catalogue}
\label{sec:catalogue}
% ============================================================================

\noindent\emph{Generated equivalent:} \code{python profiler.py -{}-catalogue}.
That output is built from the registry at runtime and cannot drift; this table
is hand-maintained. Defaults marked below are applied automatically and can be
overridden by passing the same keyword to the task. An asterisk marks the
default variant.

\begin{longtable}{@{}p{1.9cm}p{0.8cm}p{5.4cm}p{6.0cm}@{}}
\toprule
\textbf{Workload} & \textbf{Bk} & \textbf{Variants} & \textbf{Applied defaults / notes} \\
\midrule
\endhead
matmul & CPU & \code{neon}$^\ast$, \code{eigen} & \code{iterations=10};
  \code{neon} is threaded \\
matmul & GPU & \code{naive}$^\ast$, \code{tiled}, \code{fp16}, \code{vec4},
  \code{fp32\_3x16u4}, \code{fp32\_3x16u4\_db}, \code{fp32\_4x16u4},
  \code{fp32\_8x16u4}, \code{anysize}, \code{fp16\_best}, \code{best} &
  \code{ldpath="."}; the six register-blocked variants are \textbf{size=512
  only} \\
matmul & DSP & \code{scalar}, \code{hvx}$^\ast$, \code{vtcm} & separate
  binaries, not flag-selected \\
matmul & NPU & \code{fp16}$^\ast$, \code{int8} &
  \code{perf\_mode=sustained}, \code{reps=16}; \code{size}\arrow\code{-{}-n},
  \code{iterations}\arrow\code{-{}-reps} \\
\midrule
saxpy & CPU & \code{scalar}, \code{neon}$^\ast$ & \code{iterations=200};
  both threaded \\
saxpy & GPU & \code{fp32}$^\ast$ & --- \\
saxpy & DSP & \code{scalar}, \code{hvx}$^\ast$ & one binary, \code{-{}-variant};
  \code{power=lab}, \code{iterations=50} \\
saxpy & NPU & \code{fp16}$^\ast$ & \code{perf\_mode=sustained}, \code{reps=5} \\
\midrule
reduction & CPU & \code{scalar}, \code{neon}$^\ast$ & \code{iterations=200};
  both threaded \\
reduction & GPU & \code{fp32}$^\ast$ & --- \\
reduction & DSP & \code{scalar}, \code{hvx}$^\ast$ & \code{-{}-variant};
  \code{power=lab}, \code{iterations=50} \\
reduction & NPU & \code{fp16}$^\ast$ & \code{perf\_mode=sustained}, \code{reps=5} \\
\midrule
convolution & CPU & \code{scalar}, \code{neon}$^\ast$, \code{int8} &
  \code{iterations=10}; \code{neon} and \code{int8} threaded, \code{scalar} not \\
convolution & GPU & \code{fp32}$^\ast$ & --- \\
convolution & DSP & \code{scalar}, \code{hvx}$^\ast$, \code{int8} & \code{-{}-variant};
  \code{power=lab}, \code{iterations=2} \\
convolution & NPU & \code{fp16}$^\ast$ & \code{perf\_mode=sustained}, \code{reps=5} \\
\midrule
nbody & CPU & \code{neon}$^\ast$ & \code{iterations=4}; threaded \\
nbody & GPU & \code{naive}, \code{tiled}$^\ast$ & one binary,
  \code{-{}-variant}; \code{iterations=5} \\
\midrule
decode & CPU & \code{neon}$^\ast$ & --- \\
decode & GPU & \code{naive}$^\ast$, \code{gemv} & \code{ldpath="."},
  \code{wg-x=64}, \code{wg-y=8}; \code{n}\arrow\code{-{}-ncols} \\
decode & NPU & \code{fp16}$^\ast$ & \code{perf\_mode=sustained},
  \code{reps=20}; \code{n}\arrow\code{-{}-ncols} \\
\midrule
cachelat & CPU & \code{latency}$^\ast$ & \code{accesses=40000000},
  \code{pattern=random} \\
morph & DSP & \code{morph}$^\ast$ & --- \\
gpu\_memlab & GPU & \code{memlab}$^\ast$ & \code{mode=coalesce} \\
\midrule
custom & all four & --- & you supply the binary name (\S\ref{sec:custom}) \\
model & CPU, GPU, NPU & \code{qnn}$^\ast$ & dispatched through
  \code{qnn-net-run} (\S\ref{sec:models}) \\
\bottomrule
\end{longtable}

\paragraph{The GPU matmul variant trap.} \code{matmul} on the GPU defaults to
\code{naive}, which is roughly 16$\times$ slower than \code{best}. That default
is deliberate --- it keeps the optimisation ladder visible --- but it makes
cross-backend comparison meaningless if you forget it. A GPU-versus-NPU matmul
run at the default variant measures an unoptimised OpenCL kernel against a
tuned tensor engine and tells you nothing about the silicon. Name
\code{variant="best"} for any cross-engine claim.

\paragraph{The register-blocked variants accept exactly one size.}
\label{sec:fixedsize}
\code{best}, \code{fp16\_best}, \code{fp32\_3x16u4}, \code{fp32\_3x16u4\_db},
\code{fp32\_4x16u4} and \code{fp32\_8x16u4} are compiled for $N=512$ and assert
on anything else. The API checks this at construction:

\begin{lstlisting}[style=plain]
ValueError: variant 'best' is compiled for size=512 only, got size=1024.
Use an any-size variant (naive, tiled, vec4, fp16, anysize) to sweep sizes.
\end{lstlisting}

\noindent \code{anysize} is the register-blocked kernel recompiled to accept
arbitrary $N$; use it when you want the blocked shape \emph{and} a size sweep.

\paragraph{The register-tile ladder.} \code{vec4} \arrow{} \code{fp32\_3x16u4}
\arrow{} \code{fp32\_4x16u4} \arrow{} \code{fp32\_8x16u4} sweeps tile shape,
which is the knob that actually moved GFLOP/s on this GPU.
\code{fp32\_4x16u4} and \code{fp16\_best} share the 4$\times$16 tile, so that
pair isolates \emph{precision} alone; \code{best} is a different (3$\times$16)
shape. More variants exist unregistered under \code{kernels/matmul/gpu/}.

% ============================================================================
\section{Parameter reference}
\label{sec:params}
% ============================================================================

\subsection{How keywords become flags}

\begin{longtable}{@{}p{4.6cm}p{4.6cm}p{5.2cm}@{}}
\toprule
\textbf{Python keyword} & \textbf{Becomes} & \textbf{Rule} \\
\midrule
\endhead
\code{size=1024} & \code{-{}-size 1024} & scalar values are stringified \\
\code{perf\_mode="sustained"} & \code{-{}-perf-mode sustained} & underscores
  become hyphens \\
\code{H=256} & \code{-{}-H 256} & case is preserved \\
\code{validate\_only=True} & \code{-{}-validate-only} & \code{True} emits a bare
  flag with no value \\
\code{iterations=None} & (dropped) & \code{None} and \code{False} are omitted \\
\bottomrule
\end{longtable}

\noindent Values are shell-quoted with \code{shlex.quote}, so strings with
spaces are safe. Entry defaults are applied first and task parameters overwrite
them. Ordering within the command is defaults-then-overrides, which is why
\code{-{}-perf-mode sustained} appears before \code{-{}-cin} in the NPU
examples.

\subsection{Parameter renaming: \texttt{param\_map}}
\label{sec:parammap}

Three registry entries rename a parameter on its way to the CLI, because the
kernel spells it differently:

\begin{longtable}{@{}p{3.2cm}p{11.5cm}@{}}
\toprule
\textbf{Entry} & \textbf{Renames} \\
\midrule
\endhead
matmul / NPU & \code{size} \arrow{} \code{-{}-n}, \code{iterations} \arrow{}
  \code{-{}-reps} \\
decode / GPU & \code{n} \arrow{} \code{-{}-ncols} \\
decode / NPU & \code{n} \arrow{} \code{-{}-ncols} \\
\bottomrule
\end{longtable}

\noindent This exists to make the 4$\times$4 matrix uniform. Before it, NPU
matmul was the one cell where \code{iterations=N} hard-failed, so a loop over
all four backends died on one engine; and \code{size} had to be passed twice
(once as \code{n} for the kernel, once as \code{size} for the FLOP model).
Now \code{matmul(Backend.NPU, size=1024)} sends \code{-{}-n 1024} \emph{and}
computes $2 \cdot 1024^3$ FLOP, because \code{flop\_fn} reads the task's
original parameters while the rename happens only on the way to the shell.

\textbf{The corresponding hazard:} a rename that is missing or wrong does not
error. Kernels ignore unknown flags silently, so the dispatch succeeds and
quietly measures the kernel's \emph{default} shape instead of yours.

\subsection{Parameters accepted by each kernel}

\begin{longtable}{@{}p{2.3cm}p{1.1cm}p{11.2cm}@{}}
\toprule
\textbf{Workload} & \textbf{Bk} & \textbf{Accepted parameters} \\
\midrule
\endhead
matmul & CPU & \code{size}, \code{iterations}, \code{seed},
  \code{validate\_only} \\
matmul & GPU & \code{size}, \code{seed}, \code{wg\_x}, \code{wg\_y},
  \code{wg\_z}; blocked variants also take \code{inner\_reps} \\
matmul & DSP & \code{size}, \code{seed} \\
matmul & NPU & \code{size} (renamed to \code{-{}-n}), \code{iterations}
  (renamed to \code{-{}-reps}), \code{perf\_mode}, \code{seed} \\
\midrule
saxpy & CPU & \code{size}, \code{alpha}, \code{iterations}, \code{seed},
  \code{validate\_only} \\
saxpy & GPU & \code{size}, \code{alpha}, \code{seed}, \code{wg\_x},
  \code{wg\_y}, \code{wg\_z} \\
saxpy & DSP & \code{size}, \code{alpha}, \code{iterations}, \code{power},
  \code{seed}, \code{validate\_only} \\
saxpy & NPU & \code{size}, \code{alpha}, \code{reps}, \code{perf\_mode},
  \code{seed} \\
\midrule
reduction & CPU & \code{size}, \code{iterations}, \code{seed},
  \code{validate\_only} \\
reduction & GPU & \code{size}, \code{seed}, \code{wg\_x}, \code{wg\_y},
  \code{wg\_z} \\
reduction & DSP & \code{size}, \code{iterations}, \code{power}, \code{seed},
  \code{validate\_only} \\
reduction & NPU & \code{size}, \code{reps}, \code{perf\_mode}, \code{seed} \\
\midrule
convolution & CPU & \code{H}, \code{W}, \code{cin}, \code{cout}, \code{k},
  \code{iterations}, \code{seed}, \code{validate\_only} \\
convolution & GPU & \code{H}, \code{W}, \code{cin}, \code{cout}, \code{k},
  \code{seed}, \code{wg\_x} \\
convolution & DSP & \code{H}, \code{W}, \code{cin}, \code{cout}, \code{k},
  \code{iterations}, \code{power}, \code{seed} \\
convolution & NPU & \code{H}, \code{W}, \code{cin}, \code{cout}, \code{k},
  \code{reps}, \code{perf\_mode} \\
\midrule
nbody & CPU & \code{bodies}, \code{iterations}, \code{seed} \\
nbody & GPU & \code{bodies}, \code{iterations}, \code{seed} \\
\midrule
decode & CPU & \code{tokens}, \code{k}, \code{n}, \code{iterations} \\
decode & GPU & \code{tokens}, \code{k}, \code{n} (\arrow\code{-{}-ncols}),
  \code{wg\_x}, \code{wg\_y} \\
decode & NPU & \code{tokens}, \code{k}, \code{n} (\arrow\code{-{}-ncols}),
  \code{reps}, \code{perf\_mode} \\
\midrule
cachelat & CPU & \code{kb}, \code{accesses}, \code{pattern}, \code{stride},
  \code{smallpages}, \code{hugepages} (both bare flags) \\
morph & DSP & \code{op} (erode/dilate), \code{impl} (student/qhl), \code{damage}
  (salt/pepper), \code{k}, \code{input} (P5 PGM, default \code{dog.pgm}),
  \code{W}, \code{H} (synthetic image only), \code{iterations}, \code{write\_pgm} \\
gpu\_memlab & GPU & \code{mode} (coalesce/divergence/spmv/densemv), \code{stride}
  (in 4-byte floats), \code{kb} (buffer size in KiB), \code{active}, \code{density},
  \code{rows}, \code{cols}, \code{local} (workgroup size), \code{reps} \\
\bottomrule
\end{longtable}

\noindent \textbf{Do not pass \code{variant} or \code{power} as ordinary
parameters on the DSP.} \code{variant} is set for you from \code{Task.variant}
on the entries that are flag-selected, and \code{power} already defaults to
\code{lab}. Override \code{power} only when you specifically want the baseline
DCVS behaviour: \code{power="baseline"}.

\noindent \textbf{\code{cachelat}: page size and run length.} This board
runs transparent huge pages in \emph{always} mode, so the buffer already gets
2\,MB pages with no flag, and \code{-p hugepages=true} changes nothing.
\code{-p smallpages=true} forces 4\,KB pages; compare it against a run
without it to see what the TLB costs. The default \code{accesses=40000000}
finishes in about 40\,ms when the buffer fits in L1, which is too short to
trust. Set \code{-p accesses=} yourself (200000000 gives about 180\,ms on a
Prime core).

\noindent \textbf{\code{morph}: residual and images.} Every run reports
\code{residual=}. That is the number of pixels where your output differs
from the same filter applied to the undamaged photo, i.e.\ damage that
survived. Lower is better. Add \code{-p write\_pgm=true} to any run and the
profiler saves before and after pictures to \code{morph\_images/} in the
container. To open them on your laptop, run
\code{docker cp lab3:/workspace/morph\_images .} in a terminal on your machine.

\noindent \textbf{\code{gpu\_memlab}: modes, units and power runs.} One binary,
four tests, picked with \code{mode}. \code{coalesce} (the default) reads with
\code{stride} floats between neighbouring work-items and reports GB/s.
\code{divergence} lets \code{active} lanes out of every 64 do the work and
reports GFLOP/s. \code{spmv} and \code{densemv} multiply a \code{rows} $\times$
\code{cols} matrix by a vector, sparse at the given \code{density} or dense.
Units: \code{stride} counts 4-byte floats, not bytes, and \code{kb} is in KiB.
Keep \code{-p kb=65536} (64\,MiB); \code{-p kb=64} is legal but runs a 64\,KiB
buffer that fits in cache. Times are device-side OpenCL events, the minimum over
\code{reps} launches, so short runs are fine for \code{p.run()}. For
\code{p.power()} use a large \code{reps} (e.g.\ 5000): the power loop repeats the
whole command, setup and buffer upload included, and with few reps the GPU is
idle for much of the window.

\subsection{Derived FLOP and byte counts follow \emph{your} parameters}

Where a kernel prints \code{flop=} or \code{bytes=}, the printed value always
wins. Where it does not, the workload's \code{flop\_fn} / \code{bytes\_fn}
computes one from the task's parameters:

\begin{longtable}{@{}p{2.1cm}p{12.6cm}@{}}
\toprule
\textbf{Workload} & \textbf{Model} \\
\midrule
\endhead
matmul & $2n^3$ from \code{size} (default 512) \\
convolution & $2 \cdot H \cdot W \cdot c_{out} \cdot k^2 \cdot c_{in}$ from
  \code{H}, \code{W}, \code{cin}, \code{cout}, \code{k} (defaults
  $256, 256, 8, 8, 3$) \\
nbody & $20 n^2$ from \code{bodies} (default 4096), with $n$ \textbf{rounded
  down to a multiple of 64} --- both kernels do that internally, so the model
  mirrors it rather than reporting a rate computed from a body count nobody
  simulated \\
decode & FLOP $2 \cdot S \cdot k \cdot n$ from \code{tokens}, \code{k},
  \code{n}; bytes $= 4kn$, \emph{weight traffic only}. That accounting is right
  while $S \ll k$, which is the regime the workload exists to show, and stops
  being right for large $S$ \\
saxpy, reduction & no model --- these kernels print their own \code{flop=} and
  \code{bytes=} \\
\bottomrule
\end{longtable}

\noindent \textbf{The consequence is worth stating plainly.} If you change the
shape on the device without changing the corresponding Python parameter, the
timing stays correct and the derived throughput goes silently wrong. The model
has no way to see the device.

\subsection{Iteration accounting, and why it is not obvious}
\label{sec:iterations}

Kernels repeat their work internally, and they disagree about whether the time
they print covers one repetition or all of them. Getting this wrong misreports
throughput by exactly the iteration count, so the API corrects for it in two
opposite directions.

\begin{longtable}{@{}p{3.0cm}p{11.7cm}@{}}
\toprule
\textbf{Situation} & \textbf{What the API does} \\
\midrule
\endhead
\textbf{Time is a total, work model is per-iteration.} The CPU kernels report
  \code{wall\_us} across every repetition, while \code{flop\_fn} describes one
  dispatch. & Scale the work \emph{up} by the repeat count. Without this, CPU
  NEON matmul at \code{iterations=20} reported 1.37 GFLOP/s instead of
  $\approx$29. \\
\textbf{Work is a total, time is per-iteration.} The DSP saxpy, reduction and
  convolution kernels \emph{print} \code{flop=}/\code{bytes=} summed across
  every iteration, while \code{dsp\_us} is measured inside the DSP and covers
  one. & Scale the work \emph{down}, flagged by
  \code{totals\_with\_per\_iteration\_time} on those entries. Without this
  every DSP figure inflates by the iteration count --- 50$\times$ on
  saxpy/reduction, 2$\times$ on convolution, at the registry defaults. \\
\bottomrule
\end{longtable}

\noindent Only \code{wall\_us} is a total. \code{execute\_us\_mean},
\code{dsp\_us}, \code{mean\_us} and \code{duration\_ns} are already
per-dispatch --- a mean, or one timed region --- and must never be divided by
the repeat count. Doing so reported an NPU convolution as 32\,\us/iter when the
real per-inference time was 172.

\code{Result.iterations} exposes the repeat count the API used, and
\code{Result.compute\_us\_per\_iter} the per-dispatch latency.
\code{summary()} prints \code{N us/iter xM} rather than a bare time whenever
$M > 1$, so the total stays recoverable and nobody reads a ten-iteration
convolution as one slow one.


% ============================================================================
\section{What the fields mean, per backend}
\label{sec:fields}
% ============================================================================

\subsection{Selection of \texttt{compute\_us}}

The parser takes the first field present, in this order. \code{Result.compute\_key}
records which one won.

\begin{longtable}{@{}p{3.6cm}p{1.6cm}p{9.4cm}@{}}
\toprule
\textbf{Field} & \textbf{Backend} & \textbf{What it measures} \\
\midrule
\endhead
\code{execute\_us\_mean} & NPU & Mean of \code{reps} graph executions, excluding
  graph finalisation. \\
\code{dsp\_us} & DSP & Time measured on the Hexagon itself, excluding FastRPC
  marshalling. Per iteration. \\
\code{mean\_us} & varies & Mean over a kernel's own measured iterations. \\
\code{duration\_ns} & GPU & OpenCL profiling event, kernel only, excluding host
  transfers. Divided by 1000 into \us. \\
\code{wall\_us} & CPU & Host-measured wall clock around the timed region.
  \textbf{The only total.} \\
\bottomrule
\end{longtable}

\noindent Every one of these is \emph{compute only}. Comparing
\code{compute\_us} across engines compares the engines; it does not tell you
which is fastest end-to-end. A kernel that prints none of them is treated as a
contract violation and reported as a failure, not as a success with a blank
time --- see \S\ref{sec:custom}.

\subsection{Transport and overhead fields}

These live in \code{Result.fields} and are never folded into
\code{compute\_us}:

\begin{longtable}{@{}p{3.4cm}p{1.6cm}p{9.6cm}@{}}
\toprule
\textbf{Field} & \textbf{Backend} & \textbf{Meaning} \\
\midrule
\endhead
\code{transfer\_ns} & GPU & Host \arrow{} device input copy. On SAXPY this runs
  from 0.76$\times$ to 3.6$\times$ the kernel time. \\
\code{readback\_ns} & GPU & Device \arrow{} host result copy. Collapses to a few
  microseconds when the output is a scalar (reduction). \\
\code{cl\_start\_ns}, \code{cl\_end\_ns} & GPU & Raw OpenCL event timestamps. \\
\code{useful\_gbps} & GPU & Some kernels compute their own rate because only
  they know what counts as \emph{useful}. \code{gpu\_memlab} is the case in
  point: the bytes the memory system moves can differ from the bytes the kernel
  actually uses, and the host cannot see that from outside. Picked up into
  \code{Result.gbps} when nothing else supplies it. \\
\code{ns\_per\_access} & CPU & \code{cachelat}'s headline number, and the ruler
  the cache hierarchy is measured with. Promoted into \code{summary()}. \\
\code{finalize\_us} & NPU & QNN graph compilation. Roughly 10\,ms for
  convolution, 4.4\,ms for reduction, and \textbf{50--110 seconds} for SAXPY,
  which has no native primitive. \\
\code{wall\_us} & DSP, NPU & Includes FastRPC / QNN transport, so the gap
  against \code{compute\_us} is the marshalling tax (up to two orders of
  magnitude on DSP streaming kernels). \\
\code{execute\_us\_min}, \code{\_max}, \code{\_cov\_pct} & NPU & Spread across
  \code{reps}; the coefficient of variation is the cheapest stability check you
  have. \\
\code{host\_wall\_us} & all & Added by the profiler, not the kernel: the
  host-observed round trip including \code{adb} latency. Always present. For
  \code{model} tasks it is the \emph{only} timing available. \\
\code{fingerprint} & all & Numerical checksum. Compare it against a
  known-good run before trusting a timing. \\
\bottomrule
\end{longtable}

\subsection{What counters mean on a non-CPU backend}
\label{sec:accelcounters}

The \code{simpleperf} prefix is applied on \emph{every} backend, not just the
CPU, and the dispatch succeeds. \textbf{What you get back is the ARM process
that drives the accelerator, not the accelerator.} A DSP run with
\code{counters="basic"} reports the IPC of the CPU core executing the FastRPC
client while the Hexagon does the arithmetic. That is a genuinely useful
measurement --- it is how you see what the host is paying to keep an engine fed
--- but it is not the engine's IPC, and there is no way to get the engine's IPC
from here. GPU, DSP and NPU hardware counters are not exposed on this board.

\subsection{Reading a result correctly}

\begin{lstlisting}
r = p.run(saxpy(Backend.NPU, size=65536, timeout=300))

r.ok                        # did it run at all
r.fingerprint               # is the answer right
r.compute_us                # engine-only time
r.compute_key               # ...measured how? ('execute_us_mean' here)
r.iterations                # ...covering how many repetitions?
r.fields["finalize_us"]     # the compile cost you would pay in production
r.wall_us                   # what an end-to-end user would actually feel
r.raw                       # everything, unmodified, when the above disagree
\end{lstlisting}

% ============================================================================
\section{Your own kernels: \texttt{custom}}
\label{sec:custom}
% ============================================================================

\code{custom()} dispatches a binary you wrote, so adding a kernel never means
editing \code{profiler.py}.

\begin{lstlisting}
from profiler import Profiler, Backend, custom

p = Profiler()
r = p.run(custom(Backend.CPU, "my_kernel", size=1048576, core="prime"))
\end{lstlisting}

\code{binary} is a filename under
\code{/data/local/tmp/profiler-kernels/custom/}; the backend chooses the
environment (library paths, affinity) exactly as for a registered kernel. Build
and push with \code{python scripts/new\_kernel.py -{}-push <name>}.

\subsubsection*{The output contract}

Your kernel must print \textbf{one timing field}, chosen from
\code{mean\_us}, \code{wall\_us}, \code{duration\_ns}, \code{dsp\_us} or
\code{execute\_us\_mean}. This is required: a run with no time is reported as a
failure, with the contract quoted in \code{Result.error}. Leaving it as a
success produced a blank summary line and no error, which is the least helpful
thing a first custom kernel can do.

Optionally print \code{flop:}, \code{bytes:}, \code{fingerprint:} and
\code{status:}. Only \code{OK}, \code{SUCCESS}, \code{PASS}, \code{PASSED} and
\code{DONE} count as success; anything else marks the result failed. Most
shipped kernels print no status at all and rely on the exit code.

Either output style is parsed: \code{key=value} tokens anywhere on a line, or a
\code{key: value} line on its own.

\subsubsection*{Two limitations}

\begin{itemize}[leftmargin=0.6cm]
  \item \code{custom()} accepts \code{name}, \code{timeout}, \code{depends\_on}
    and \code{core}, but \textbf{not} \code{threads}, \code{counters} or
    \code{per\_core}. To use those, construct the \code{Task} directly with
    \code{binary=}.
  \item \code{custom} resolves binaries under \code{custom/} only. A registered
    workload's binary cannot be driven through it --- which is why
    \code{cachelat} is in the registry rather than left to \code{custom}: it
    needs counters and lives in \code{cachelat/}.
\end{itemize}

% ============================================================================
\section{Converted networks: \texttt{model}}
\label{sec:models}
% ============================================================================

Networks are not kernels. They are DLC files executed by QAIRT's own
\code{qnn-net-run}, and the engine is chosen by which QNN backend library is
loaded --- \code{libQnnCpu.so}, \code{libQnnGpu.so} or \code{libQnnHtp.so}. One
DLC runs on any of the three, so \textbf{the engine is the variable and the
model is fixed}, which is the whole point of the case study.

\begin{lstlisting}
from profiler import Profiler, Backend, model

p = Profiler()
r = p.run(model(Backend.NPU, "minilm", reps=10))
print(r.wall_us)                       # the only timing qnn-net-run gives

for b in (Backend.CPU, Backend.GPU, Backend.NPU):
    print(p.time_model(b, "mobilenetv2", reps=100))
\end{lstlisting}

\noindent The DSP is deliberately absent: QAIRT ships a QNN DSP backend, but
provisioning does not push it, so listing it would advertise a dispatch that
cannot run.

\subsection{The model registry}

\code{MODELS} maps a short name to its directory, DLC, input list, MAC count
and kind. All eleven live under \code{/data/local/tmp}, staged separately from
the kernel tree.

\begin{longtable}{@{}p{2.7cm}p{1.5cm}p{1.6cm}p{8.2cm}@{}}
\toprule
\textbf{Name} & \textbf{Kind} & \textbf{MACs} & \textbf{Notes} \\
\midrule
\endhead
\code{mobilenetv2} & cnn & 301\,M & depthwise-separable, linear topology \\
\code{mnv2\_s1} & cnn & 120.8\,M & stage 1 of 3 (40\% of MACs) \\
\code{mnv2\_s2} & cnn & 104.6\,M & stage 2 of 3 (35\%) \\
\code{mnv2\_s3} & cnn & 75.3\,M & stage 3 of 3 (25\%) \\
\code{minilm} & encoder & 15\,M & 6-layer BERT encoder, seq=128 \\
\code{distilgpt2} & decoder & 88.4\,M & 6-layer decoder \textsc{prefill},
  seq=128, no KV cache \\
\code{qwen05\_dec32} & decoder & 494\,M & Qwen2.5-0.5B decode, 2.52\,GB DLC \\
\code{llama1b\_dec32} & decoder & 1236\,M & Llama-3.2-1B decode, 6.00\,GB DLC \\
\code{sd\_text\_encoder} & clip & 6.7\,G & CLIP text encoder, once per image \\
\code{sd\_unet} & unet & 68\,G & distilled SD UNet, \textbf{once per denoising
  step} \\
\code{sd\_vae\_decoder} & vae & 1.3\,T & VAE decoder; the most engine-sensitive
  stage on this board (8$\times$ spread) \\
\bottomrule
\end{longtable}

\noindent The three \code{mnv2\_s*} stages exist so a network can be cut and the
pieces placed on different engines. The diffusion trio is the pipeline where
that actually pays: its stage boundary is a 64\,KB latent, cheap enough for a
genuine multi-engine pipeline, which the LLM's 4.7\,MB KV handoff never allowed.

\textbf{The decode models are fp32 and that is not a choice.} The CPU backend
refuses fp16 DLCs (Discovery~71), so the DLCs are large and decode is
weight-bound. \textbf{The diffusion text encoder takes token
\emph{embeddings}, not ids}: the converter folds integer inputs away
(Discovery~77), so the 77-row lookup happens on the host and only the twelve
transformer layers run on device.

\subsection{\texttt{model(backend, name, **params)}}

\begin{longtable}{@{}p{2.3cm}p{12.4cm}@{}}
\toprule
\textbf{Argument} & \textbf{Meaning} \\
\midrule
\endhead
\code{backend} & \code{Backend.CPU}, \code{.GPU} or \code{.NPU}. Anything else
  raises. \\
\code{name} & A key of \code{MODELS}. Unknown names raise, listing the known
  ones. \\
\code{reps} & \code{-{}-num\_inferences}: how many inferences one dispatch
  performs. Default 1. \\
\code{output\_dir} & Where \code{qnn-net-run} writes outputs on device.
  Defaults to \code{run\_<model>\_<backend>}, and is deleted before each run. \\
\code{name=} & Task name, defaulting to \code{model.<name>.<backend>}. \\
\bottomrule
\end{longtable}

\noindent A \code{model} result has \textbf{no \code{compute\_us}} ---
\code{qnn-net-run} prints no timing field the parser recognises. Its
\code{wall\_us} is filled in by the profiler from the host round trip, and
\code{ok} additionally requires the string \emph{``Finished Executing''} in the
output. A model that was never staged fails with a message naming the command
that stages it, rather than a bare \code{cd: No such file or directory}.

\subsection{\texttt{time\_model(backend, name, reps=100)}}

Graph preparation is paid once per dispatch and differs enormously by backend
--- 0.16\,s on CPU, 0.52\,s on GPU, 0.59\,s on HTP for MobileNetV2 --- so a
single timed run measures mostly compilation. \code{time\_model} runs the model
twice, once with \code{reps=1} and once with \code{reps+1}, and subtracts. The
subtraction cancels preparation and the constant \code{adb} round trip.

\begin{lstlisting}
{'model': 'mobilenetv2', 'backend': 'npu', 'ok': True,
 'per_inference_ms': 0.812, 'prepare_plus_one_ms': 604.3,
 'inferences_per_s': 1231.5, 'reps': 100, 'macs': 301000000.0}
\end{lstlisting}

\noindent \textbf{It refuses rather than returning a negative number.} The
subtraction only cancels preparation if preparation costs the same in both
runs, and it often does not --- a cold DLC read from storage can dominate the
first run, after which the \code{reps+1} run comes back \emph{faster} in wall
time and the per-inference figure goes negative. Reporting $-0.75$\,ms as a
measurement is worse than reporting nothing, so the call returns
\code{ok: False} with an error saying to raise \code{reps} until the
per-inference signal outweighs the one-off load.

% ============================================================================
\section{Power estimation: \texttt{power()}}
\label{sec:power}
% ============================================================================

\begin{lstlisting}
p.power(matmul(Backend.CPU, size=1024))
# {'available': True, 'power_w': 6.59, 'energy_j': 230.6, 'window_s': 35,
#  'engine': 'cpu', 'boundary': '12V supply input',
#  'scope': 'whole_board_not_per_engine', 'board_match': 'calibration_board'}
\end{lstlisting}

\subsection{Three things this is not}

Each matters more than the number itself.

\begin{description}[leftmargin=1.2cm,style=nextline,font=\normalfont\itshape]
  \item[Not a measurement.] It is a calibrated \emph{model} of 12\,V
    supply-input power, fitted against a bench supply. This board has no
    readable PMIC rail --- the battery node is a \code{Debug\_Board} dummy
    (Discovery~57) --- so there is nothing to read. The returned
    \code{quality} is \code{modeled\_ballpark} for that reason.
  \item[Not per-engine.] It predicts what the \emph{whole board} draws.
    ``GPU watts'' is not a thing this can produce.
  \item[Not per-dispatch.] The model is validated only over intervals of at
    least \code{minimum\_validated\_interval\_s} (30\,s), so the task is looped
    to fill \code{seconds} and the answer describes that window. A single
    200\,\us{} dispatch has no modeled power. Asking for fewer seconds than the
    artifact's validated interval returns \code{available: False}.
\end{description}

\subsection{Arguments}

\begin{longtable}{@{}p{2.3cm}p{12.4cm}@{}}
\toprule
\textbf{Argument} & \textbf{Meaning} \\
\midrule
\endhead
\code{task} & The \code{Task} to loop. One task, one engine. \\
\code{seconds} & Window length, default 35. Must be $\geq$ the artifact's
  validated interval (30\,s). \\
\code{artifact} & Path to a power-model JSON. Default: the best available for
  the attached board --- the SM8850 CPU/GPU artifact first, then the SM8750
  ones, so a Hexagon calibration appears automatically the moment one lands. \\
\bottomrule
\end{longtable}

\subsection{What each engine returns}

\begin{longtable}{@{}p{1.6cm}p{13.1cm}@{}}
\toprule
\textbf{Engine} & \textbf{Status} \\
\midrule
\endhead
CPU & Fitted from measurements. \\
GPU & Fitted from measurements. \\
NPU & \emph{Declared}, not fitted: 5\,W fp16, 4\,W int8, hand-entered from
  operator readings. Returned with \code{validated: False} so a reader can tell
  a fitted number from a typed-in one, and returned \emph{immediately} without
  running a window --- running a 35\,s NPU window and then ignoring its features
  would be theatre, and it is exactly the run that drives the die to 98\,\textdegree C
  (Discovery~73). \\
DSP & \code{available: False}, with a reason. An unsigned PD cannot cast the
  HAP power vote the calibration needs. \textbf{Reporting unavailable is the
  correct answer, not a fallback number} --- see ADRs 0019, 0020 and 0022. \\
\bottomrule
\end{longtable}

\subsection{Board transfer}

The model's coefficients were fitted on one specific board. \code{board\_match}
says which case you are in:

\begin{itemize}[leftmargin=0.6cm]
  \item \code{calibration\_board} --- fitted on this exact unit.
  \item \code{transferred\_eas\_verified} --- fitted elsewhere, and this board's
    kernel energy model was checked against the calibration board's before the
    coefficients were applied. More uncertainty in the absolute watts;
    \emph{none} in a comparison between two runs on this board.
\end{itemize}

\noindent A serial mismatch with no EAS verification blocks the estimate
outright. Two same-model HDK8850 units differed by 4.44\%/4.72\% in their
vendor EAS tables, so the silicon genuinely is not identical.

\subsection{Why there is no multi-engine power call}

The model's concurrent-path terms are not summable
(\code{concurrent\_path\_terms\_summable: false}), so there is no call that
gives the power of a multi-engine pipeline. Measure each stage and add the
\emph{energies}.

\subsection{Operational notes}

The call takes about a minute and prints a progress line to stderr, so it does
not look hung. The device-side loop is given \code{seconds + 25} so it cannot
end mid-window, which means it outlives the measurement; \code{power()} reaps
it afterwards, because otherwise a second \code{power()} call lands while the
first loop is still running and is refused with \emph{``compute clients are
already running''}.

\code{reason} explains why there is no number; \code{warnings} says what to
distrust about one that is present. Read both.

% ============================================================================
\section{Command-line interface}
% ============================================================================

The module is runnable directly, for single points, for the catalogues, and for
the demo tour.

\begin{lstlisting}[style=sh]
python profiler.py --demo
python profiler.py --catalogue
python profiler.py --list-counters

python profiler.py matmul npu
python profiler.py matmul gpu --variant best -p size=512
python profiler.py reduction cpu --core prime -p size=4194304
python profiler.py convolution cpu --core gold --threads 6 --counters cache --per-core
python profiler.py convolution npu -p cin=64 -p cout=64 -v
python profiler.py cachelat cpu -p kb=256 -p pattern=random --repeat 5
python profiler.py custom cpu --binary my_kernel -p size=1048576
python profiler.py saxpy dsp --serial 5cabf4fd
\end{lstlisting}

\begin{longtable}{@{}p{3.5cm}p{11.2cm}@{}}
\toprule
\textbf{Flag} & \textbf{Meaning} \\
\midrule
\endhead
\code{-{}-demo} & Run the live tour. Also the default when no workload is
  given. \\
\code{-{}-catalogue} & Print every workload, backend, variant and applied
  default, generated from the registry. \\
\code{-{}-list-counters} & Print the presets, the 92 verified events by group,
  and the 43 blocked ones with reasons. Exits without touching the device. \\
\code{-{}-serial SERIAL} & Target one board when several are attached. \\
\code{-v}, \code{-{}-verbose} & Print each generated device command. \\
\code{-{}-variant NAME} & Kernel variant. \\
\code{-{}-core SPEC} & Affinity: a core (\code{7}), a mask (\code{0xc0}), or
  \code{prime}/\code{gold}/\code{all}. \\
\code{-{}-threads N} & Worker threads (threaded variants only). \\
\code{-{}-counters SPEC} & A preset name or a comma-separated event list. \\
\code{-{}-per-core} & Break counter totals down per CPU. \\
\code{-{}-counter-cpus SPEC} & Count per-CPU system-wide on these CPUs
  (e.g.\ \code{0-5}) instead of per-process. \\
\code{-p KEY=VAL} & A kernel parameter; repeatable. Values are coerced:
  integers, then floats, then \code{true}/\code{false} into bare flags, else
  left as a string. \\
\code{-{}-repeat N} & Run the whole dispatch $N$ times and report mean, range
  and coefficient of variation. \\
\code{-{}-binary NAME} & With workload \code{custom}: the binary under the
  device \code{custom/} directory. Rejected with any other workload. \\
\bottomrule
\end{longtable}

\subsection{\texttt{-{}-repeat}: process repeats, not kernel iterations}

\code{-{}-repeat N} runs $N$ separate \emph{processes}. That is the point:
process-level repeats re-pay startup and let the governor settle differently
each time, so the spread they expose is the spread a real caller would see. A
kernel's own \code{-{}-iterations} cannot show you that --- it amortises it
away.

The summary picks whichever headline metric the workload actually reports
(GFLOP/s, else GB/s, else \code{wall\_us}):

\begin{lstlisting}[style=plain]
    1  matmul.cpu    105234.1 us/iter x10    29.3 GFLOP/s  fp=1.3587e+03
    2  matmul.cpu    106881.7 us/iter x10    28.9 GFLOP/s  fp=1.3587e+03
    3  matmul.cpu    105902.4 us/iter x10    29.1 GFLOP/s  fp=1.3587e+03

  n=3  GFLOP/s: mean 29.1  range 28.9-29.3  CoV 0.7%
\end{lstlisting}

\noindent Failed runs are excluded from the statistics and counted in a note.
If the runs produce more than one distinct fingerprint, a warning says so ---
they did not all compute the same answer, and the timing spread is the least of
your problems.

\noindent \textbf{One run is not a measurement.} \code{-{}-repeat} exists to
make that easy to act on.

% ============================================================================
\section{Recipes}
% ============================================================================

\subsection{One roofline point on every engine}

\begin{lstlisting}
from profiler import Profiler, Backend, convolution

p = Profiler()
for r in p.run_serial([convolution(b, cin=8, cout=8) for b in Backend]):
    if r.ok:
        print(f"{r.task.backend}  {r.compute_us:10.1f} us  {r.gflops:8.1f} GFLOP/s")
    else:
        print(f"{r.task.backend}  FAILED: {r.error}")
\end{lstlisting}

\subsection{A parameter sweep written to CSV}

\begin{lstlisting}
import csv
from profiler import Profiler, Backend, matmul

p = Profiler()
with open("matmul_gpu.csv", "w", newline="") as fh:
    w = csv.writer(fh)
    w.writerow(["size", "variant", "compute_us", "gflops", "fingerprint"])
    for size in (128, 256, 512, 1024):
        for variant in ("naive", "tiled", "anysize"):   # any-size variants only
            r = p.run(matmul(Backend.GPU, size=size, variant=variant))
            if r.ok:
                w.writerow([size, variant, r.compute_us, r.gflops, r.fingerprint])
\end{lstlisting}

\subsection{Contention: the same batch, both ways}

\begin{lstlisting}
batch = [matmul(Backend.CPU, size=512), matmul(Backend.GPU, size=512, variant="best")]

clean     = p.run(batch, parallel=False)
contended = p.run(batch, parallel=True)

for a, b in zip(clean, contended):
    print(f"{a.task.name:20s} {a.compute_us:9.1f} -> {b.compute_us:9.1f} us "
          f"({b.compute_us / a.compute_us:.2f}x)")
\end{lstlisting}

\subsection{CPU cluster and core study, with counters}

\begin{lstlisting}
for spec in ("prime", "gold", 7, 0):
    r = p.run(reduction(Backend.CPU, size=4 << 20, core=spec,
                        counters="basic", name=f"core_{spec}"))
    d = r.derived
    print(f"{str(spec):6s} {r.compute_us:9.1f} us  "
          f"{d.get('effective_ghz', 0):.2f} GHz  "
          f"{d.get('clock_of_max_pct', 0):5.1f}% of max  IPC={d.get('ipc', 0):.2f}")
\end{lstlisting}

\noindent Name the cluster. Under the default mask \code{clock\_of\_max\_pct}
is absent, because cores 4--7 straddle the boundary (\S\ref{sec:affinity}).

\subsection{Thread scaling, done so the answer means something}

\begin{lstlisting}
for n in (1, 2, 3, 4, 5, 6):
    r = p.run(convolution(Backend.CPU, cin=32, cout=32, iterations=20,
                          threads=n, core="gold", variant="neon",
                          counters="cache", name=f"conv_{n}t"))
    print(f"{n} threads  {r.compute_us_per_iter:9.1f} us/iter  "
          f"IPC={r.derived.get('ipc', 0):.2f}  "
          f"L2 miss={r.derived.get('l2d_miss_pct', 0):.2f}%")
\end{lstlisting}

\noindent Three things make this measurement honest rather than misleading:
\code{iterations=20} so thread spawn is amortised (\S\ref{sec:threads});
\code{core="gold"} so six threads have six cores; and \code{variant="neon"}
because \code{scalar} would raise.

\subsection{Scalar versus vectorised, on one engine}

\begin{lstlisting}
for variant in ("scalar", "neon"):
    r = p.run(saxpy(Backend.CPU, size=1 << 22, variant=variant, core="prime"))
    print(f"{variant:7s} {r.gbps:6.1f} GB/s")
\end{lstlisting}

\subsection{A pipeline expressed as a DAG}

\begin{lstlisting}
results = p.run_dag([
    reduction(Backend.CPU, size=1 << 20, name="stage1"),
    convolution(Backend.GPU, cin=16, cout=16, name="stage2a", depends_on=["stage1"]),
    convolution(Backend.NPU, cin=16, cout=16, name="stage2b", depends_on=["stage1"]),
    matmul(Backend.CPU, size=512, name="stage3", depends_on=["stage2a", "stage2b"]),
])
for name in ("stage1", "stage2a", "stage2b", "stage3"):
    print(results[name].summary())
\end{lstlisting}

\subsection{Comparing engines on one network}

\begin{lstlisting}
for b in (Backend.CPU, Backend.GPU, Backend.NPU):
    t = p.time_model(b, "mobilenetv2", reps=100)
    if t["ok"]:
        print(f"{b}  {t['per_inference_ms']:7.3f} ms  "
              f"{t['inferences_per_s']:8.1f} inf/s  "
              f"(prepare {t['prepare_plus_one_ms']:.0f} ms)")
    else:
        print(f"{b}  {t['error']}")
\end{lstlisting}

\subsection{Checking what will run, without a board}

\begin{lstlisting}
from profiler import build_command, DEVICE_BASE, WORKLOADS, Backend, matmul

for variant in WORKLOADS["matmul"].kernels[Backend.GPU].variants:
    print(f"{variant:16s} {build_command(matmul(Backend.GPU, variant=variant), DEVICE_BASE)}")
\end{lstlisting}

\subsection{A guarded sweep that reports its own failures}

\begin{lstlisting}
results, failures = [], []
for size in (256, 512, 1024, 2048):
    r = p.run(matmul(Backend.DSP, size=size, timeout=600))
    (results if r.ok else failures).append(r)

print(f"{len(results)} ok, {len(failures)} failed")
for r in failures:
    print(f"  {r.task.name} size={r.task.params['size']}: {r.error}")

fps = {round(r.fingerprint, 6) for r in results if r.fingerprint is not None}
if len(fps) > 1:
    print(f"WARNING: {len(fps)} distinct fingerprints -- not all the same answer")
\end{lstlisting}

% ============================================================================
\section{Extending the wrapper}
% ============================================================================

\subsection{Running a kernel without touching this file}

Use \code{custom()} (\S\ref{sec:custom}). Nothing needs registering. This is the
right answer for a one-off, for coursework, and for anything you are still
iterating on.

\subsection{Adding a variant to an existing kernel}

Push the binary, then add one entry to the \code{variants} dict:

\begin{lstlisting}
Backend.GPU: KernelEntry("gpu", {"naive": "matmul_opencl_fp32",
                                 "mine":  "matmul_opencl_fp32_mine"},
                         "naive", ldpath="."),
\end{lstlisting}

\noindent If it is compiled for one shape only, add it to \code{fixed\_size} so
the guard fires on the host. If it accepts \code{-{}-threads}, add it to
\code{threaded\_variants}.

\subsection{Adding a backend to an existing workload}

Add a \code{KernelEntry} under the new \code{Backend} key. The environment
(working directory, library paths, affinity, counter prefix) is applied by
backend, so nothing else changes. If the kernel spells a parameter differently,
add a \code{param\_map} entry rather than making callers remember it.

\subsection{Adding a whole workload}

Add a \code{Workload} to \code{WORKLOADS} with its \code{kernels} table, and a
\code{flop\_fn} / \code{bytes\_fn} if the kernels do not print their own counts.
Then add a module-level constructor:

\begin{lstlisting}
def myworkload(backend: Backend, **params) -> Task:
    """One line saying what the workload is for and where it sits on the roofline."""
    return _task("myworkload", backend, params)
\end{lstlisting}

\noindent Add the name to \code{\_\_all\_\_}, or \code{from profiler import *}
will silently not bind it.

\subsection{Adding a new timing field}

Append it to \code{\_COMPUTE\_FIELDS} in priority order. If the new field is a
\emph{total} across repetitions rather than a per-dispatch figure, also add it
to \code{\_TOTAL\_COMPUTE\_FIELDS} --- otherwise throughput is wrong by exactly
the iteration count (\S\ref{sec:iterations}).

\subsection{Adding a derived counter metric}

Append a tuple to \code{\_DERIVED}: \code{(name, numerator candidates,
denominator candidates, scale)}. The first candidate present wins, and the
metric is computed only when both sides exist --- so a metric that needs events
the caller did not select simply does not appear, rather than appearing as
zero.

% ============================================================================
\section{Semantics, caveats, and measurement validity}
\label{sec:caveats}
% ============================================================================

\subsection{The caveats that change conclusions}

\begin{description}[leftmargin=0.8cm,style=nextline,font=\normalfont\bfseries]
  \item[NPU and DSP do not run in parallel.] They share the Hexagon cDSP.
    Parallel submission looks right from the host and serialises in hardware.
    \S\ref{sec:parallel}.
  \item[The DSP cannot pin its own clock.]
    \label{sec:dspclock}
    HAP power voting is refused from the unsigned PD we run in
    (\code{0x80000414}), with no testsig installed. Compute over FastRPC works
    fine; only the power vote does not. \textbf{Consequence: DSP numbers are
    unpinned and vary more than the other backends.} Two identical DSP runs
    differing by 50\% is normal, not a bug. Use \code{-{}-repeat} and quote a
    spread.
  \item[A busy NPU makes the DSP faster.] Same shared DVFS domain. A DSP kernel
    measured beside an NPU peer can be 1.59$\times$ faster than the same kernel
    alone. Every DSP figure here is the idle-Hexagon case.
  \item[The default affinity mask spans two clusters.] \code{core=None} is cores
    4--7: two Gold, two Prime. Name the cluster for anything you will quote.
    \S\ref{sec:affinity}.
  \item[Defaults are not the fastest variant.] GPU matmul defaults to
    \code{naive}, about 16$\times$ off \code{best}. Deliberate, and fatal to a
    cross-engine comparison if forgotten.
  \item[NPU SAXPY needs a long timeout.] QNN has no native SAXPY operator, so
    the graph decomposes into thousands of tiles and \code{finalize\_us} runs
    50--110 seconds. The default 180\,s timeout is not always enough; pass
    \code{timeout=300}. The wait \emph{is} the finding.
  \item[GPU counters are not available.] Adreno perfcounters are blocked on this
    build. Counters on a GPU task measure the host process.
    \S\ref{sec:accelcounters}.
  \item[Never run \code{adb root}.] It restarts \code{adbd} in a domain that
    cannot reach the OpenCL ICD, so every GPU dispatch fails with
    \code{clGetPlatformIDs} $= -1001$. Counters use \code{su 0}, which is a
    different and safe mechanism. \S\ref{sec:su}.
  \item[The NPU can destabilise the board in still air.] A sustained fp16 graph
    reached 98\,\textdegree C with mitigation engaged, at only 4--5\,W, and the
    board dropped off USB. Power density, not total power. Do not leave a
    sustained NPU sweep running unattended without airflow.
  \item[Thermals are tighter than the previous board.] This HDK cannot hold its
    top GPU or CPU operating points under load in still air. Several published
    calibration setpoints sit below the top of the ladder for that reason.
\end{description}

\subsection{Checking a number before you quote it}

\begin{enumerate}[leftmargin=0.8cm]
  \item \code{r.ok} is \code{True}.
  \item \code{r.fingerprint} matches a known-good run. A fast wrong answer is
    still wrong.
  \item \code{r.compute\_key} is the field you think it is, and
    \code{r.iterations} is what you asked for.
  \item \code{r.below\_max\_clock} is \code{False}, or you know why it is not.
  \item \code{r.multiplexed} is \code{False} if you are reading counter values.
  \item You ran it more than once (\code{-{}-repeat}) and the coefficient of
    variation is small enough to support the claim.
  \item The comparison is like-for-like: same variant class, same sizes, same
    serial-versus-parallel mode.
\end{enumerate}

% ============================================================================
\section{Troubleshooting}
% ============================================================================

\begin{longtable}{@{}p{5.1cm}p{9.6cm}@{}}
\toprule
\textbf{Symptom} & \textbf{Cause and fix} \\
\midrule
\endhead
\code{ValueError: unknown workload} & Typo, or a workload that does not exist.
  The message lists the known ones. \\
\code{ValueError: X has no Y kernel} & That workload does not run on that
  engine --- \code{nbody} has no NPU, \code{cachelat} is CPU-only. The message
  lists the backends it does have. \\
\code{ValueError: ... has no variant} & Typo. The message lists every valid
  variant and often a ``did you mean''. \\
\code{ValueError: compiled for size=512 only} & A register-blocked GPU variant.
  Use \code{anysize}, \code{naive}, \code{tiled}, \code{vec4} or \code{fp16}
  to sweep sizes. \S\ref{sec:fixedsize}. \\
\code{ValueError: ... does not accept -{}-threads} & That variant is
  single-threaded; the message names the threaded ones. \S\ref{sec:threads}. \\
\code{ValueError: counter ... reads a constant 0} & The event is not
  implemented on this silicon. The message names the substitute.
  \S\ref{sec:counters}. \\
Result \code{ok=False}, error mentions \code{No such file} & The binary was
  never pushed. Re-run \code{provision\_device.py}, then
  \code{adb shell ls} the directory to confirm --- \emph{``1 file pushed''} is
  not evidence. \\
Every counter run fails, no counters at all & \code{su 0} is unavailable, or
  you are on a build where the shell domain cannot open the PMU and \code{su}
  is absent. Everything except counters still works. \\
Exit code 124 in \code{error} & Timeout. Raise \code{timeout=}; for NPU SAXPY
  use 300\,s or more. \\
Blank summary, \code{ok=False}, ``printed no timing field'' & A custom kernel
  that does not meet the output contract. \S\ref{sec:custom}. \\
Throughput looks wrong by exactly $10\times$ or $50\times$ & An iteration
  accounting mismatch. Check \code{r.compute\_key} and \code{r.iterations}.
  \S\ref{sec:iterations}. \\
GFLOP/s right, timing right, but for the wrong shape & A parameter the kernel
  ignored. Turn on \code{verbose=True} and read the command. \\
GPU dispatches all fail with \code{-1001} & Somebody ran \code{adb root}.
  \code{adb unroot} and reconnect. \\
\code{power()} returns \code{available: False} for the DSP & Working as
  designed --- the DSP is not calibrated and will not be guessed at.
  \S\ref{sec:power}. \\
\code{power()} says ``compute clients are already running'' & A previous
  \code{power()} loop has not been reaped. Wait 25\,s, or
  \code{adb shell pkill -f profiler\_power\_loop}. \\
\code{time\_model} returns negative per-inference time & Load noise exceeded
  the signal. Raise \code{reps}. \S\ref{sec:timemodel}. \\
Results differ between two identical DSP runs & Expected: the DSP clock is
  unpinned. \S\ref{sec:dspclock}. \\
\code{adb} dies mid-sweep & Intermittent on this board. The provisioner is
  idempotent and sweeps are restartable; re-run. \\
\bottomrule
\end{longtable}

% ============================================================================
\appendix
\section{Counter catalogue}
\label{sec:counterappendix}
% ============================================================================

\noindent\emph{Generated equivalent:} \code{python profiler.py -{}-list-counters}.
92 events are verified on this silicon, across ten groups. Each was run against
a memory-heavy workload on both clusters and returned a non-zero count.
\textbf{Note this checks that a
counter counts, not that its value is correct.}

\begin{longtable}{@{}p{2.2cm}p{0.9cm}p{11.4cm}@{}}
\toprule
\textbf{Group} & \textbf{N} & \textbf{What it covers} \\
\midrule
\endhead
\code{core} & 3 & \code{cpu-cycles}, \code{instructions}, \code{bus-cycles}. \\
\code{branch} & 9 & Retired and speculative branches, predictor hits and
  misses, split by direct / indirect / return. \\
\code{cache} & 20 & L1D, L1I and L2 accesses, refills, reads, writes,
  writebacks, allocations, long-latency misses, and L1D refills split by inner
  (on-cluster) vs outer (off-cluster). \\
\code{prefetch} & 7 & Accesses generated by the \emph{hardware} prefetcher
  (\code{*-hwprf}) versus by software \code{PRFM} instructions
  (\code{*-prfm}). A sequential stream drives \code{hwprf} hard; a pointer
  chase does not, and that contrast is the experiment. \\
\code{tlb} & 18 & L1 dTLB/iTLB, L2 unified TLB, refills, and page-table walks
  including walk depth (\code{*-step}) and whether the walk resolved to a large
  or a small page. \\
\code{stalls} & 9 & Frontend and backend stalls in both cycles and issue slots,
  plus memory-attributed backend stalls. Slots are finer than cycles. \\
\code{mix} & 8 & Retired vs speculative instructions and micro-ops, loads,
  stores, integer data-processing, and FP/SIMD. \textbf{NEON counts under
  \code{raw-vfp-spec}}, not \code{raw-ase-inst-spec}, which reads zero. \\
\code{memory} & 4 & Memory accesses (total, read, write) and bus/interconnect
  accesses. \\
\code{exceptions} & 2 & Exceptions taken and returned. \\
\code{aliases} & 12 & simpleperf's friendlier \code{hw-cache} names for events
  already listed above (\code{L1-dcache-loads}, \code{dTLB-load-misses}, and so
  on). Same silicon events. \\
\bottomrule
\end{longtable}

\subsection{Derived metrics}

Computed automatically from whichever events are present. A metric whose
numerator or denominator is missing simply does not appear.

\begin{longtable}{@{}p{4.3cm}p{10.4cm}@{}}
\toprule
\textbf{Metric} & \textbf{Definition} \\
\midrule
\endhead
\code{ipc}, \code{cpi} & instructions / cycles, and its reciprocal \\
\code{branch\_miss\_pct} & mispredicts / branches \\
\code{l1d\_miss\_pct}, \code{l1i\_miss\_pct}, \code{l2d\_miss\_pct} &
  refills / accesses, per level \\
\code{l1dtlb\_miss\_pct}, \code{l2dtlb\_miss\_pct} & TLB refills / TLB
  accesses \\
\code{pagewalks\_per\_kinst} & data page-table walks per thousand instructions \\
\code{frontend\_stall\_pct}, \code{backend\_stall\_pct} & stall cycles / total
  cycles \\
\code{mem\_stall\_pct} & memory-attributed backend stalls / cycles \\
\code{slot\_backend\_pct} & backend-lost issue slots / total lost slots \\
\code{l1d\_wb\_per\_kinst} & L1D writebacks per thousand instructions \\
\code{l1d\_hwprf\_pct}, \code{l2d\_hwprf\_pct} & share of cache traffic the
  \emph{hardware prefetcher} generated rather than the program \\
\code{spec\_waste\_pct} & $1 - \text{retired}/\text{speculative}$: the fraction
  of executed instructions that never retired \\
\code{effective\_ghz} & measured clock, from cycles over simpleperf's window
  (\S\ref{sec:clock}) \\
\code{clock\_of\_max\_pct} & effective clock as a percentage of the cluster
  ceiling \\
\code{scaling\_cur\_freq\_ghz} & the governor's \emph{requested} clock, sampled
  during the run \\
\bottomrule
\end{longtable}

\section{Glossary of the numbers this document quotes}

\begin{longtable}{@{}p{4.3cm}p{10.4cm}@{}}
\toprule
\textbf{Constant} & \textbf{Value and meaning} \\
\midrule
\endhead
\code{DEVICE\_BASE} & \code{/data/local/tmp/profiler-kernels} \\
\code{MODEL\_DIR} & \code{/data/local/tmp} --- networks are staged outside the
  kernel tree \\
\code{CLUSTER\_MAX\_GHZ} & \code{\{"gold": 3.6288, "prime": 4.3968\}} --- the
  top of each cluster's published OPP ladder \\
\code{CLOCK\_FLAG\_RATIO} & \code{0.80} --- below this fraction of the ceiling,
  \code{below\_max\_clock} is set \\
\code{MAX\_HW\_COUNTERS} & \code{12} --- PMU counters per core; beyond this
  simpleperf multiplexes \\
default \code{timeout} & 180.0 seconds \\
default \code{core} & mask \code{f0}, cores 4--7 \\
default \code{power()} window & 35 s, against a 30 s validated minimum \\
\bottomrule
\end{longtable}

\vfill
\noindent\rule{\linewidth}{0.4pt}\\[2pt]
\small
Document revision 2.0. The catalogues here are hand-maintained;
\code{-{}-catalogue} and \code{-{}-list-counters} are generated and
authoritative.

\end{document}
